\documentclass[9pt,letterpaper]{extarticle}

\usepackage[letterpaper,
            left=0.70in,right=0.35in,top=0.42in,bottom=0.45in,
            headheight=10pt,headsep=8pt]{geometry}
\usepackage{amsmath,amssymb,mathtools,bm}
\allowdisplaybreaks
\usepackage{booktabs,array,longtable,tabularx,multirow}
\usepackage{needspace}
\usepackage{enumitem}
\usepackage{listings}
\usepackage{algpseudocode}
\usepackage[compact]{titlesec}
\usepackage{microtype}
\usepackage[table]{xcolor}
\usepackage[T1]{fontenc}
\usepackage{lmodern,anyfontsize}
\usepackage[colorlinks=true,linkcolor=blue!50!black,urlcolor=blue!50!black]{hyperref}
\usepackage{fancyhdr}
\usepackage{claimledger}
\usepackage{f2zspec}

\definecolor{specnavy}{HTML}{17365D}
\definecolor{specblue}{HTML}{315B7D}
\definecolor{specmuted}{HTML}{667085}
\definecolor{specpale}{HTML}{F3F4F6}

\setlength{\parindent}{0pt}
\linespread{0.96}
\setlength{\parskip}{1pt}
\setlength{\abovedisplayskip}{2pt}
\setlength{\belowdisplayskip}{2pt}
\setlength{\jot}{1pt}
\setlength{\emergencystretch}{2em}
\tolerance=1700
\setlist{nosep,leftmargin=1.35em}
\titlespacing*{\section}{0pt}{0.48\baselineskip}{0.18\baselineskip}
\titlespacing*{\subsection}{0pt}{0.30\baselineskip}{0.08\baselineskip}

\pagestyle{fancy}
\fancyhf{}
\fancyhead[L]{\footnotesize\sffamily\color{specmuted}F2Z PCS general specification}
\fancyhead[R]{\footnotesize\sffamily\color{specmuted}chunked mod-$q$ / recursive Ligerito}
\fancyfoot[C]{\thepage}
\renewcommand{\headrulewidth}{0.25pt}

\renewcommand{\clgboardtitle}{Reduction ledger}

\lstdefinestyle{wire}{
  basicstyle=\ttfamily\scriptsize,
  columns=fullflexible,
  keepspaces=true,
  frame=single,
  rulecolor=\color{specmuted!35},
  backgroundcolor=\color{specpale},
  xleftmargin=0.35em,
  framexleftmargin=0.35em,
  aboveskip=3pt,
  belowskip=3pt
}

\newenvironment{normative}
  {\par\smallskip\begingroup
   \setlength{\leftskip}{0.7em}\setlength{\rightskip}{0.7em}%
  %  \color{specnavy}\noindent\textbf{Normative v1 proof system.}
   }
  {\par\endgroup\smallskip}

\begin{document}
\clgreset

\begin{center}
  {\Large\bfseries\color{specnavy}The F2Z polynomial commitment scheme}\\[2pt]
  {\small Relations, reductions, and a normative prover--verifier specification}
\end{center}

\begin{normative}
F2Z is a polynomial-commitment layer for proving the Lin-BitsDual relation.
Given a public vector $\bm v$ and a value $\mu\in\Fq$, the relation asserts
\[
  \left\langle\bm v,\pi_q(\bm h)\right\rangle_{\Fq}=\mu,
  \qquad
  \pi_2(\bm h)=\bm M\pi_2(\bm f),
  \qquad
  O=\Enc_{\mathcal C}\!\left(\Pack(\pi_2(\bm f))\right),
\]
where $\bm f$ is the committed binary witness, $\bm h$ is the active witness
used by the application, and $\bm M$ is a public developer-selected
virtualization map.
Here $\pi_2:\mathbb Z\to\Ftwo$ denotes coordinatewise reduction modulo $2$.
The application PIOP supplies this linear claim to F2Z.  Under the configured
layout of $\bm h$, an admissible claim has a separable coefficient matrix
$V[r,c]=a_r a'_c$; the public factors $(\bm a,\bm a')$ are part of the F2Z
instance.

Drawing inspiration from binary-field techniques used in systems such as Flock and Binius64, F2Z combines the following techniques:

\begin{itemize}[leftmargin=1.4em]
 \item \textbf{Binary packing.}
       The committed witness bits are packed into elements of the binary
       extension field $\K$ and encoded using the binary linear code
       $\mathcal C$. This allows the commitment and opening protocols to
       operate on packed field elements.

 \item \textbf{Witness virtualization.}
       The public map $\bm M$ relates the committed witness $\bm f$ to the
       witness $\bm h$ on another ring/field. This permits constraints to be layered on top of each other while making the witness $\bm f$ to be small.

 \item \textbf{Chunked integer lifting.}
       Coefficients in $\Fq$ are lifted to canonical integers and split into
       bounded chunks. This preserves the original modular claim while keeping
       every intermediate fold within the injective range of the binary-field
       exponent encoding.

 \item \textbf{Grand-product encoding.}
       Because the witness coordinates are binary, each bounded linear fold can
       be expressed as a product identity over $\K$. This converts the original
       linear claim into product claims suited to binary-field verification.

 \item \textbf{GKR and LinCheck reductions.}
       GKR verifies the product computation succinctly, while the row LinCheck
       reduces the resulting weighted claim to a multilinear-extension
       evaluation of the active witness.

 \item \textbf{Adjoint pullback.}
       The verifier transports the final coefficient through
       $\bm M^{\mathsf T}$, converting a claim on $\bm h$ into a claim on the
       originally committed witness $\bm f$ without another commitment or
       protocol message.

 \item \textbf{Conditional reduction, ring switching, and batching.}
       The coefficient LinCheck is used only when the pulled-back coefficient
       is not already in the required MLE form. A ring switch then converts the
       resulting evaluations into packed linear claims, which are batched into
       a single recursive Ligerito opening.
\end{itemize}

Together, these techniques allow F2Z to prove inner products claims over $\Fq$, and
therefore support constraints beyond $\Ftwo$, while retaining a binary-field
commitment. This specification fixes the Lin-BitsDual relation, its
prover--verifier reductions, and its concrete transcript and encoding rules.
\end{normative}

\section{Overview and interfaces}
\label{sec:overview}

\begingroup\small

\subsection{Committed and virtual witness views}
\label{sec:public-objects}

The PCS receives a committed binary witness $\bm f$.  Before proving, the
developer fixes the matrix layouts used by the commitment and the PIOP.
The same configuration may optionally include a public virtualization map:

\begin{itemize}[leftmargin=1.4em]
 \item \textbf{Direct.}
       The virtualization map is omitted and the PIOP uses the
       committed witness directly, so $\bm h=\bm f$.
 \item \textbf{Virtualized.}
       The PIOP uses a binary witness $\bm h$ related to the
       committed witness by a nontrivial developer-selected public matrix
       $\bm M$.
\end{itemize}

Throughout, $\bm f$ is the commitment-facing vector and $\bm h$ is the
PIOP-facing vector.  When $\bm M$ is omitted, every uniform formula
below interprets it as the identity map.

The virtual witness is
\begin{equation}
 \boxed{\quad
 \pi_2(\bm h)=\bm M\pi_2(\bm f).
 \quad}
 \label{eq:virtual-witness}
\end{equation}
Here $\pi_2:\mathbb Z\to\Ftwo$, $x\mapsto x\bmod 2$, acts coordinatewise.
Since $\bm f$ and $\bm h$ are binary, Equation~\eqref{eq:virtual-witness}
uniquely determines $\bm h$.  If virtualization is enabled, the developer
fixes $\bm M$ before proving; it is not chosen by the prover.

Only $\bm f$ is committed.  The developer-configured matrices $F$ and $H$ are
reshapings of $\bm f$ and $\bm h$, not additional witnesses.  The commitment is
\begin{equation}
 \bm f\in\bitsset^{n_f},\qquad
 \operatorname{vec}(F):=\pi_2(\bm f),\qquad
 P:=\Pack(F),\qquad
 O:=\Enc_{\mathcal C}(P),\qquad
 \rho:=\Root(O).
 \label{eq:overview-commitment}
\end{equation}
Neither $H$, the layout of $\bm h$, nor any value derived from $H$ is encoded
in $O$.

\paragraph{Adjoint convention.}
For the standard inner product over a field $\mathbb E$, the adjoint of a
matrix $\bm T$ is its transpose $\bm T^{\mathsf T}$:
\begin{equation}
 \langle\bm u,\bm T\bm v\rangle_{\mathbb E}
 =\langle\bm T^{\mathsf T}\bm u,\bm v\rangle_{\mathbb E}.
 \label{eq:adjoint-convention}
\end{equation}
Throughout this specification, applying the adjoint means using this identity
to move a public matrix across the inner product.  It is a deterministic
rewrite and sends no protocol message.

\subsection{PIOP linear-claim handoff}
\label{sec:piop-handoff}

The end-to-end proof begins with a constraint-satisfaction relation over $\Fq$
evaluated on the virtual witness $\pi_q(\bm h)$.  The PIOP reduces that relation
to a general Lin-BitsDual claim consumed by F2Z.  The input Sumcheck phase of
$\Pi_1$ shapes that claim before the integer-fold phase:
\begin{equation}
 \begin{aligned}
  (x;\pi_q(\bm h))\in\mathcal R_{\rm CS}
  &\xRightarrow{\ \Pi_{\rm PIOP}\ }
    \left\langle\bm v,\pi_q(\bm h)\right\rangle_{\Fq}=\mu \\
  &\xRightarrow{\ \Pi_1\ }
    \mathcal R_{\rm PESAT}[\bm h] \\
  &\xRightarrow{\ \mathrm{F2Z}\ }
    \text{opening against $\rho$}.
 \end{aligned}
 \label{eq:piop-f2z-flow}
\end{equation}
The PIOP is defined by the constraint system and is not part of F2Z itself.
The end-to-end interaction includes it as the subprotocol that produces the
F2Z handoff.  F2Z receives the arbitrary public coefficient $\bm v$ and does
not inspect the equations or reductions that produced it.  The checked input
Sumcheck is internal to $\Pi_1$; it does
not factor $\bm v$, but replaces the original scalar claim by an
evaluation-shaped intermediate claim before producing F2Z PESAT instances.

\endgroup

\section{Instance and parameters}
\label{sec:instance}

\begingroup\small

Section~\ref{sec:instance} is the canonical registry for persistent and
cross-phase F2Z notation.  A definition gives the exact value denoted by a
symbol, its type gives the ambient set or interface class, and its rule records
the normative constraint on its use.  Section~\ref{sec:encoding} specifies how
the transcript derives the non-input symbols registered here.

\newcommand{\DefTableHead}{%
  \toprule
  \cellcolor{specnavy}\textcolor{white}{\textbf{Symbol}} &
  \cellcolor{specblue}\textcolor{white}{\textbf{Definition}} &
  \cellcolor{spectableteal}\textcolor{white}{\textbf{Type}} &
  \cellcolor{specnavy}\textcolor{white}{\textbf{Rule}} \\
  \midrule}

\newcommand{\DefTableSetup}{%
  \setlength{\tabcolsep}{3pt}%
  \renewcommand{\arraystretch}{1.10}%
  \arrayrulecolor{spectableline}%
  \rowcolors{2}{white}{spectablerow}}

\subsection{Algebraic domains and canonical maps}
\label{sec:domains}

\begingroup
\scriptsize
\DefTableSetup
\begin{tabularx}{\linewidth}{@{}>{\color{specnavy}\raggedright\arraybackslash}p{0.78in}
 >{\raggedright\arraybackslash}X
 >{\color{spectableteal}\raggedright\arraybackslash}p{1.25in}
 >{\raggedright\arraybackslash}X@{}}
\DefTableHead
$\mathbb Z$ & Exact integer domain & Ring \john{don't need} &\\
\midrule
$\Ftwo$ & Binary field & $\mathbb F_2$ &\\
\midrule
$\Fq$ & $q=2^{100}-15$ & Finite field &\\
\midrule
$\K$ & $\Ftwo[X]/(X^{128}+X^7+X^2+X+1)$ & $\mathbb F_{2^{128}}$ & Used for packing, products, GKR, and openings. \\
\midrule
$\pi_2$ & $x\mapsto x\bmod2$ & $\mathbb Z\to\Ftwo$ & \\
\midrule
$\operatorname{lift}_2(b)$ & The unique $u\in\{0,1\}\subset\mathbb Z$ satisfying $\pi_2(u)=b$ & $\Ftwo\to\{0,1\}$ & Canonical bit representative; not an inverse of $\pi_2$ on all of $\mathbb Z$. \\
\midrule
$\pi_q$ & $x\mapsto x\bmod q$ & $\mathbb Z\to\Fq$ & \\
\midrule
$\operatorname{lift}_q$ & $\operatorname{lift}_q(a)$ is the unique integer in $[0,q)$ satisfying $\pi_q(\operatorname{lift}_q(a))=a$ & $\Fq\to[0,q)$ & Canonical integer lift; no alternative representative is accepted. \\
\midrule
$\iota_{2\to\K}$ & Canonical subfield embedding & $\Ftwo\hookrightarrow\K$ & Applied coordinatewise when a binary witness enters a $\K$-claim. \\
\midrule
$\beta_v$ & $X^v$ in the fixed polynomial basis of $\K$ & Element of $\K$ & $0\le v<128$, ordered by increasing $v$. \\
\bottomrule
\end{tabularx}
\endgroup

\subsection{Dimensions and index conventions}
\label{sec:dimensions}

\begingroup
\scriptsize
\DefTableSetup
\begin{tabularx}{\linewidth}{@{}>{\color{specnavy}\raggedright\arraybackslash}p{0.78in}
 >{\raggedright\arraybackslash}X
 >{\color{spectableteal}\raggedright\arraybackslash}p{1.25in}
 >{\raggedright\arraybackslash}X@{}}
\DefTableHead
$n_f$ & Length of the committed witness & $\mathbb N$ & $n_f=2^{m_f}$ and $128\mid n_f$. \\
\midrule
$n_h$ & Length of the virtual witness & $\mathbb N$ & $n_h=2^{m_h}$. \\
\midrule
$m_f$ & Boolean dimension of $\bm f$ & $\mathbb N$ & \\
\midrule
$m_h$ & Boolean dimension of $\bm h$ & $\mathbb N$ & \\
\midrule
$d_{F,r},d_{F,c}$ & Row and column dimensions of $F$ & $\mathbb N^2$ & $d_{F,r}+d_{F,c}=m_f$. \\
\midrule
$d_r,d_c$ & Row and column dimensions of $H$ & $\mathbb N^2$ & $d_r+d_c=m_h$ and $2^{d_r}q<2^{127}$. \\
\midrule
$m_p$ & Packed-message dimension & $\mathbb N$ & $m_p=m_f-7$. \\
\midrule
$r$ & Row index of $H$ & $0\le r<2^{d_r}$ &  \\
\midrule
$c$ & Column index of $H$ & $0\le c<2^{d_c}$ & \\
\midrule
$j$ & Packed-element index & $0\le j<n_f/128$ & \\
\midrule
$v$ & Basis-coordinate index in $K$ & $0\le v<128$ & index to reference a coefficient for a packed galios number \\
\bottomrule
\end{tabularx}
\endgroup

\subsection{Witnesses, layouts, and virtualization}
\label{sec:witness-layouts}

\begingroup
\scriptsize
\DefTableSetup
\begin{tabularx}{\linewidth}{@{}>{\color{specnavy}\raggedright\arraybackslash}p{0.78in}
 >{\raggedright\arraybackslash}X
 >{\color{spectableteal}\raggedright\arraybackslash}p{1.25in}
 >{\raggedright\arraybackslash}X@{}}
\DefTableHead
$\bm f$ & Commitment-facing binary witness & $\bitsset^{n_f}$ & The PCS commits only to this vector. \\
\midrule
$F$ & Developer-configured reshaping of $\pi_2(\bm f)$ & $\Ftwo^{2^{d_{F,r}}\times2^{d_{F,c}}}$ & $\operatorname{vec}(F)=\pi_2(\bm f)$. The developer reshapes it in terms of $2^{d_{F,r}}\times2^{d_{F,c}}$. \\
\midrule
$\bm h$ & PIOP-facing binary witness & $\bitsset^{n_h}$ & Equals $\bm f$ when virtualization is omitted. \\
\midrule
$H$ & Developer-configured reshaping of $\pi_2(\bm h)$ & $\Ftwo^{2^{d_r}\times2^{d_c}}$ & $\operatorname{vec}(H)=\pi_2(\bm h)$. \\
\midrule
$H_{\mathbb Z}$ & Entrywise $\operatorname{lift}_2(H)$ & $\{0,1\}^{2^{d_r}\times2^{d_c}}\subset\mathbb Z^{2^{d_r}\times2^{d_c}}$ & Used in exact integer folds. \\
\midrule
$H_q$ & Entrywise $\pi_q(H_{\mathbb Z})$ & $\Fq^{2^{d_r}\times2^{d_c}}$ & Used in the Lin-BitsDual claim and the input Sumcheck. \\
\midrule
$H_{\K}$ & Entrywise $\iota_{2\to\K}(H)$ & $\K^{2^{d_r}\times2^{d_c}}$ & Used in PESAT, GKR, and MLE claims. \\
\midrule
$\operatorname{vec}$ & Flattening of a configured matrix layout & Matrix-to-vector map & The developer fixes its coordinate order before proving. \\
\midrule
$\bm M$ & Effective virtualization map & $\Ftwo^{n_h\times n_f}$ & Supplied by the developer when virtualization is enabled.  Otherwise $n_h=n_f$, both layouts preserve the same flattened coordinate order, $\bm h=\bm f$, and $\bm M:=I$. \\
\midrule
$\bm M_{\K}$ & Scalar extension of the effective $\bm M$ & $\K^{n_h\times n_f}$ & Its transpose is the adjoint from Equation~\eqref{eq:adjoint-convention}. \\
\midrule
$\bm f_{\K},\bm h_{\K}$ & Coordinatewise images under $\iota_{2\to\K}$ & $\K^{n_f}$ and $\K^{n_h}$ & Used only in $\K$-linear claims. \\
\bottomrule
\end{tabularx}
\endgroup

The matrices $F$ and $H$ are views of the witness vectors, not additional
witnesses.  Their shapes and coordinate orders are developer configuration
choices fixed before proving.  The configured Boolean address of $H$ is
$z=(r,c)\in\bitsset^{d_r}\times\bitsset^{d_c}$; its vectorization and every
multilinear extension in this specification use this same row/column bit
split.  In particular, reshaping a vector $\bm u\in\Fq^{n_h}$ into a matrix
$U$ means $\operatorname{vec}(U)=\bm u$ under that configured order, not an
independently chosen row-major convention.

\subsection{Packing and commitment}
\label{sec:packing-commitment}

\begingroup
\scriptsize
\DefTableSetup
\begin{tabularx}{\linewidth}{@{}>{\color{specnavy}\raggedright\arraybackslash}p{0.78in}
 >{\raggedright\arraybackslash}X
 >{\color{spectableteal}\raggedright\arraybackslash}p{1.25in}
 >{\raggedright\arraybackslash}X@{}}
\DefTableHead
$\Pack$ & Packs consecutive groups of $128$ coordinates into $\K$ & $\Ftwo^{n_f}\to\K^{n_f/128}$ & Uses the ordered basis $(\beta_v)_{v<128}$. \\
\midrule
$P$ & $P[j]=\sum_{v=0}^{127}\operatorname{vec}(F)[128j+v]\beta_v$ & $\K^{n_f/128}$ & Derived only from the configured view $F$. \\
\midrule
$F_v$ & $F_v[j]=\operatorname{vec}(F)[128j+v]$ & $\Ftwo^{2^{m_p}}$ & The $v$-th bit slice of the packed message, so $P[j]=\sum_vF_v[j]\beta_v$. \\
\midrule
$\mathcal C$ & Configured encoding code & Public code &  \\
\midrule
$\Enc_{\mathcal C}$ & Canonical encoder for $\mathcal C$ & $\K^{n_f/128}\to\K^{n_O}$ & \\
\midrule
$O$ & $\Enc_{\mathcal C}(P)$ & $\K^{n_O}$ & Encodes $P$ and never the virtual view $H$. \\
\midrule
$\Root$ & Configured Merkle-root algorithm & $\K^{n_O}\to\bitsset^{\lambda}$ & Constructs a Merkle tree on $\Enc_{\mathcal C}(P)$. \\
\midrule
$\rho$ & $\Root(O)$ & $\bitsset^{\lambda}$ & Public commitment to $O$. \\
\bottomrule
\end{tabularx}
\endgroup

\subsection{Structured ring-switch coefficients}
\label{sec:ring-switch-coefficients}

\begingroup
\scriptsize
\DefTableSetup
\begin{tabularx}{\linewidth}{@{}>{\color{specnavy}\raggedright\arraybackslash}p{0.78in}
 >{\raggedright\arraybackslash}X
 >{\color{spectableteal}\raggedright\arraybackslash}p{1.25in}
 >{\raggedright\arraybackslash}X@{}}
\DefTableHead
$A_f[v,j]$ & $\bm a_f[128j+v]$ in the configured commitment order & $\K^{128\times2^{m_p}}$ & Pack-boundary view of the public coefficient produced by the adjoint. \\
\midrule
$\Theta$ & $\{0,\ldots,t_{\rm sw}-1\}$ in canonical order & Finite nonempty index set & Derived by the developer configuration from $\bm a_f$, with $t_{\rm sw}<2^{32}$. \\
\midrule
$A_\theta,B_\theta$ & $A_\theta[v]\in\K$ and $B_\theta[j]\in\K$ & $\K^{128}$ and $\K^{2^{m_p}}$ & Must satisfy $A_f[v,j]=\sum_{\theta\in\Theta}A_\theta[v]B_\theta[j]$ for every $(v,j)$. \\
\midrule
$\mathsf{RingDecomp}$ & $\bm a_f\mapsto(\Theta,(A_\theta,B_\theta)_{\theta\in\Theta})$ & Public deterministic map & The verifier derives the same ordered decomposition; no proof field is used. \\
\bottomrule
\end{tabularx}
\endgroup

Every coefficient admits the canonical fallback
$\Theta=\{0,\ldots,127\}$,
$A_\theta[v]=1$ when $v=\theta$ and $0$ otherwise, and
$B_\theta[j]=A_f[\theta,j]$.  A developer configuration may instead provide a
smaller structured decomposition, such as the translated-equality terms
induced by rotations, shifts, or other public witness wiring.

\subsection{Lin-BitsDual public claim}
\label{sec:linbits-instance}

\begingroup
\scriptsize
\DefTableSetup
\begin{tabularx}{\linewidth}{@{}>{\color{specnavy}\raggedright\arraybackslash}p{0.78in}
 >{\raggedright\arraybackslash}X
 >{\color{spectableteal}\raggedright\arraybackslash}p{1.25in}
 >{\raggedright\arraybackslash}X@{}}
\DefTableHead
$\bm v$ & Public coefficient vector & $\Fq^{n_h}$ & Uses the configured coordinate order of $H$. \\
\midrule
$V$ & $\operatorname{vec}(V)=\bm v$ in the configured order of $H$ & $\Fq^{2^{d_r}\times2^{d_c}}$ & Arbitrary; no rank or separability condition is imposed. \\
\midrule
$\mu$ & Claimed inner product $\langle\bm v,\pi_q(\bm h)\rangle_{\Fq}=\langle\bm v,\operatorname{vec}(H_q)\rangle_{\Fq}$ & $\Fq$ & Initial target of the input Sumcheck. \\
\midrule
$\mathcal I_{\rm LinBitsDual}$ & $(\bm v,\mu)$ together with the developer configuration & Public F2Z instance & Binds the arbitrary input claim, both layouts, and the effective virtualization map. \\
\midrule
$\bm s_h=(\bm s_r,\bm s_c)$ & Input-Sumcheck terminal point & $\Fq^{m_h}$ & Transcript-derived with $|\bm s_r|=d_r$ and $|\bm s_c|=d_c$. \\
\midrule
$\kappa$ & $\widetilde{\bm v}(\bm s_h)$ & $\Fq$ & Computed by the verifier; never supplied by the prover. \\
\midrule
$\tau$ & Input-Sumcheck terminal value & $\Fq$ & Reduced target satisfying $\kappa\widetilde H_q(\bm s_h)=\tau$. \\
\bottomrule
\end{tabularx}
\endgroup

After the input Sumcheck, the verifier locally derives
$a_r:=\eq_{\Fq}(r,\bm s_r)$ and
$a'_c:=\kappa\eq_{\Fq}(c,\bm s_c)$.  Consequently,
\begin{equation}
 W[r,c]:=\kappa\eq_{\Fq}((r,c),\bm s_h)
        =\eq_{\Fq}(r,\bm s_r)\,
         \kappa\eq_{\Fq}(c,\bm s_c)=a_ra'_c,
 \label{eq:derived-core-coefficient}
\end{equation}
only the new coefficient $W$ is separable.  Equation~\eqref{eq:derived-core-coefficient}
is a consequence of the equality polynomial at the Sumcheck terminal point;
it is not a decomposition of the arbitrary input matrix $V$.

\subsection{Integer lifts and fold parameters}
\label{sec:core-parameters}

\begingroup
\scriptsize
\DefTableSetup
\begin{tabularx}{\linewidth}{@{}>{\color{specnavy}\raggedright\arraybackslash}p{0.86in}
 >{\raggedright\arraybackslash}X
 >{\color{spectableteal}\raggedright\arraybackslash}p{1.27in}
 >{\raggedright\arraybackslash}X@{}}
\DefTableHead
$\gamma_r$ & $\operatorname{lift}_q(a_r)$ & $[0,q)\subset\mathbb Z$ & Unique canonical integer lift of $a_r$. \\
\midrule
$\widehat\mu_c$ & $\sum_{r<2^{d_r}}\gamma_rH_{\mathbb Z}[r,c]$ & $[0,2^{127})\subset\mathbb Z$ & The bound follows from $2^{d_r}q<2^{127}$ and is required for exponent injectivity. \\
\midrule
$\mu_c$ & Prover-supplied claimed fold & $[0,2^{127})\subset\mathbb Z$ & Accepted only if its exponent matches PESAT and $\tau=\sum_ca'_c\pi_q(\mu_c)$. \\
\midrule
$g$ & Product generator & $\K^\times$ & $\operatorname{ord}(g)=2^{128}-1$, making exponentiation injective on accepted folds. \\
\midrule
$\nu_c$ & $g^{\mu_c}$ & $\K^\times$ & Must equal $\prod_r\bigl(1+(g^{\gamma_r}-1)H_{\K}[r,c]\bigr)$. \\
\bottomrule
\end{tabularx}
\endgroup

\subsection{Protocol interfaces}
\label{sec:protocol-interfaces}

\begingroup
\scriptsize
\DefTableSetup
\begin{tabularx}{\linewidth}{@{}>{\color{specnavy}\raggedright\arraybackslash}p{0.9in}
 >{\raggedright\arraybackslash}X
 >{\raggedright\arraybackslash}X@{}}
\toprule
\cellcolor{specnavy}\textcolor{white}{\textbf{Symbol}} &
\cellcolor{specblue}\textcolor{white}{\textbf{Definition}} &
\cellcolor{specnavy}\textcolor{white}{\textbf{Rule}} \\
\midrule
$\Pi_1$ & $\begin{gathered}\langle\bm v,\pi_q(\bm h)\rangle_{\Fq}=\mu\\[-2pt] \Longrightarrow\quad g^{\mu_c}=\displaystyle\prod_r\left(1+(g^{\gamma_r}-1)H_{\K}[r,c]\right)\end{gathered}$ & Runs the input Sumcheck, derives the canonical integer weights, and checks the fold bounds, exponentiation, and reconstruction of the reduced scalar claim. \\
\midrule
$\Pi_{\rm GKR}$ & $\begin{gathered}g^{\mu_c}=\displaystyle\prod_r\left(1+(g^{\gamma_r}-1)H_{\K}[r,c]\right)\\[-2pt] \Longrightarrow\ \displaystyle\sum_z\omega_{\xi,z}H_{\K}[z]=\sigma\end{gathered}$ & \\
\midrule
$\Pi_{\rm MLE}$ & $\begin{gathered}\displaystyle\sum_z\omega_{\xi,z}H_{\K}[z]=\sigma\\[-2pt] \Longrightarrow\ \widetilde H_{\K}(\bm r_h)=\eta\end{gathered}$ & Typically reduced by $\mathsf{LinCheck}_H$. \\
\midrule
$\operatorname{Adj}_{\bm M}$ & $\langle\bm u,\bm M_{\K}\bm v\rangle_{\K}\longmapsto\langle\bm M_{\K}^{\mathsf T}\bm u,\bm v\rangle_{\K}$ & \\
\midrule
$\Pi_{\rm iopp}$ & $\begin{gathered}\langle\bm a_f,\bm f_{\K}\rangle_{\K}=\eta\\[-2pt] \Longrightarrow\ \langle B,P\rangle_{\K}=\beta\ \Longrightarrow\ \mathsf{accept}\end{gathered}$ & Performs one ring switch and one recursive Ligerito opening. \\
\bottomrule
\end{tabularx}
\endgroup

\endgroup

\section{Relations, reductions, and proof objects}
\label{sec:encoding}

\begingroup\footnotesize

\subsection{How to read this section}

A paper-level encoded relation contains triples $(x,O;w)$: $x$ is the explicit
instance, $O$ is the implicit encoded oracle, and $w$ is the witness.  For a public
witness map $T$, the binding condition is
\[
 O=\Enc_{\mathcal C}(T(w)).
\]
Merkle compilation commits to $O$ with the public root
$\rho:=\Root(O)$; $\rho$ is not a proof field.
Every relation below is additionally parameterized by the fixed public
context $\chi:=(\mathsf{pp},t,s,W,g)$ and its derived dimensions.  The context
is omitted from displayed instances and is carried unchanged by every edge.

The edge notation is
\[
 \mathcal R_0\equiv_F\mathcal R_1
 \quad\text{deterministic retyping},\qquad
 \mathcal R_0\xRightarrow{\Pi}\mathcal R_1
 \quad\text{interactive reduction},
\]
\[
 \mathcal R_0\mathrel{\dashrightarrow}\mathcal R_1
 \quad\text{conditional extension or classification only}.
\]
Each subsection defines one source relation and lists its direct outgoing
reductions.  A proof fragment lists transmitted values only.  Challenges,
successor coefficients, exponentiated values, batching coefficients, and
work tables are verifier-derived.

\subsection{The application PIMSAT relations}

Let $S:=R_M[X]$, and let $M\subseteq S$ be a statement-bound finite free
$R_M$-submodule.  Fix a public instance $\bm x\in S^{\ell_x}$, a bounded
witness domain $\mathcal W_B\subseteq M^r$, ideals $I_a\lhd S$, and constraint
polynomials $p_j^{(a)}\in S[\bm X,\bm Y]$.  The paper first uses the
unencoded relation
\begin{align}
 \RPIMSAT^{\rm app}:=\Bigl\{(\bm x,\bot;\bm f):\;&
 \bm f\in\mathcal W_B,
 \quad p_j^{(a)}(\bm x,\bm f)\in I_a\quad(\forall a,j)\Bigr\}.
 \label{eq:rel-pimsat-unencoded}
\end{align}
Encoding the witness gives the encoded application relation
\begin{align}
 \RPIMSAT^{\rm enc}:=\Bigl\{(\bm x,O;\bm f):\;&
 \bm f\in\mathcal W_B,
 \quad O=\Enc_{\mathcal C}\!\left(
          \Pack(\pi_2(\bits(\bm f)))\right),\notag\\[-2pt]
 &p_j^{(a)}(\bm x,\bm f)\in I_a\quad(\forall a,j)\Bigr\}.
 \label{eq:rel-pimsat-app}
\end{align}

\paragraph{Direct reductions.}
\begin{itemize}

\item \textbf{Reduction 0: Compile the witness oracle.}
\[
 \RPIMSAT^{\rm app}
 \xRightarrow{\mathsf{Encode}}
 \RPIMSAT^{\rm enc}.
\]
The IOP supplies the encoded oracle $O$; Merkle compilation exposes
$\rho=\Root(O)$.  The constraint predicates are unchanged.

\item \textbf{Reduction 1: Reduce ideal membership to application PESAT.}
\[
 \RPIMSAT^{\rm enc}
 \xRightarrow{\Pi_{{\rm app},1}}
 \RPESAT^{\rm app}.
\]
\begin{itemize}
 \item \textbf{Proof:} $\pi_{{\rm app},1}$ carries the PIOP messages that turn
 ideal-membership constraints into field equations.
 \item \textbf{Result:} the verifier derives $K_{\rm app}$,
 $\psi_{\rm app}$, $\widehat{\bm x}$, and $\widehat{\bm p}$; $O$ is unchanged.
 \item \textbf{Scope:} application-supplied, before the F2Z PCS core.
\end{itemize}

\item \textbf{Classification: Equality and congruence constraints.}
$I_a=(0)$ gives equality; nonzero $I_a$ gives congruence or quotient-ring
constraints.  Both use the same reduction.

\end{itemize}

\subsection{The encoded application PESAT relation}

Fix the verifier-derived application field $K_{\rm app}$, a public map
$\psi_{\rm app}:M\to K_{\rm app}$, and projected constraint polynomials
$\widehat p_i\in K_{\rm app}[\bm X,\bm Y]$.  The encoded application PESAT
relation retains the original bounded witness:
\begin{align}
 \RPESAT^{\rm app}:=
 \Bigl\{((\widehat{\bm x},\widehat{\bm p}),O;\bm f):\;&
 \bm f\in\mathcal W_B,
 \quad O=\Enc_{\mathcal C}\!\left(
             \Pack(\pi_2(\bits(\bm f)))\right),\notag\\[-2pt]
 &\widehat p_i(\widehat{\bm x},\psi_{\rm app}(\bm f))=0
   \quad(\forall i)\Bigr\}.
 \label{eq:rel-pesat-app}
\end{align}

\paragraph{Direct reductions.}
\begin{itemize}

\item \textbf{Reduction 1: Linearize the application PESAT claim.}
\[
 \RPESAT^{\rm app}
 \xRightarrow{\Pi_{{\rm app},2}}
 \RLMFB.
\]
\begin{itemize}
 \item \textbf{Proof:} $\pi_{{\rm app},2}$ carries the application reduction
 from the projected constraints to one scalar linear claim.
 \item \textbf{Result:} the verifier derives $(\bm v,\mu)$ and $\psi$; the
 same witness and $O$ pass to $\RLMFB$.
 \item \textbf{Scope:} application-supplied.  Together,
 $\pi_{\rm PIOP}:=(\pi_{{\rm app},1},\pi_{{\rm app},2})$ is the complete
 application proof fragment.
\end{itemize}

\item \textbf{Classification: A single linear constraint is LMFB.}
If $\widehat{\bm p}$ is one linear polynomial, the relation is already
Lin-ModFieldBin; no extra message is needed.

\end{itemize}

\subsection{The Lin-ModFieldBin relation}

Let $(\gamma_j)_{j<D_M}$ be the fixed basis of $M$, and assume a
statement-bound bijection
$\bits:\mathcal W_B\leftrightarrow\bitsset^n$ with fixed reconstruction:
\[
 f_i=\sum_{k<B}\sum_{j<D_M}b_{i,k,j}2^k\gamma_j,
 \qquad b_{i,k,j}\in\bitsset,
 \qquad \bm b=\bits(\bm f).
\]
At this seam a higher-rank target has already been expanded into scalar
coordinates; write $\psi:M\to R$ for the resulting $\varphi$-semilinear map.
\Needspace{7\baselineskip}
The module relation is
\begin{align}
 \RLMFB:=\Bigl\{((\bm v,\mu),O;\bm f):\;&
 \bm f\in\mathcal W_B,
 \quad\bm v\in R^r,
 \quad\mu\in R,
 \quad O=\Enc_{\mathcal C}\!\left(
                 \Pack(\pi_2(\bits(\bm f)))\right),\notag\\[-2pt]
 &\sum_i v_i\psi(f_i)=\mu\quad\text{in }R\Bigr\}.
 \label{eq:rel-lmfb}
\end{align}

\paragraph{Direct reductions.}
\begin{itemize}

\item \textbf{Reduction 1: Deterministic Bitify.}

\[
 \RLMFB
 \equiv_{\mathsf{Bitify}}
 \RLBD,
 \qquad
 u_{i,k,j}:=
 v_i\,\varphi(\iota_{R_M}(2^k))\,\psi(\gamma_j).
\]
The witness is retyped as $\bm b:=\bits(\bm f)$; the same encoded oracle $O$
remains bound to it.  Bitify transmits nothing and performs no new Booleanity
test.

\item \textbf{Required preprocessing for a higher-rank target.}
If $N$ has rank greater than one, first expand it into scalar LMFB claims.
Then bitify each coordinate or use an application-defined batch; no canonical
batching proof is specified here.

\end{itemize}

\subsection{The Lin-BitsDual relation}

The core bit relation is
\begin{align}
 \RLBD:=\Bigl\{((\bm u,\mu),O;\bm b):\;&
 \bm b\in\bitsset^n,
 \quad\bm u\in R^n,
 \quad\mu\in R,
 \quad O=\Enc_{\mathcal C}(\Pack(\pi_2(\bm b))),\notag\\[-2pt]
 &\langle\iota_R(\bm b),\bm u\rangle_R=\mu\Bigr\}.
 \label{eq:rel-lbd}
\end{align}

Suppose the public coefficient has the paper's tensor form
$\bm u=\bm u^{(1)}\otimes\bm u^{(2)}$, and reshape the committed bits as
$\Bbit[c,i]$.  With canonical integer representatives
$\gamma_i$ of $u_i^{(1)}$, the prover exposes column values
\[
 \mu_c:=\sum_i\gamma_i\Bbit[c,i]\in\Z.
\]
The verifier immediately checks the tensor reconstruction
\begin{equation}
 \sum_c u_c^{(2)}\,\iota_R(\mu_c)=\mu.
 \label{eq:pi1-paper-reconstruction}
\end{equation}
It then derives $y_i:=g^{\gamma_i}$ and $\nu_c:=g^{\mu_c}$, obtaining
\[
 \prod_i\bigl(1+(y_i-1)\Bbit[c,i]\bigr)=\nu_c
 \qquad(\forall c).
\]

\paragraph{Direct reductions.}
\begin{itemize}

\item \textbf{Reduction 1: Produce the F2Z encoded PESAT claim.}
\[
 \RLBD
 \xRightarrow{\Pi_1}
 \RPESAT(\K,Q_{\Ftwo\!\Z},\mathcal C,\Pack).
\]
\begin{itemize}
 \item \textbf{Proof:} the paper sends $\pi_1=(\mu_c)_c$.  The verifier checks
 Equation~\eqref{eq:pi1-paper-reconstruction} before accepting the successor.
 \item \textbf{Result:} the verifier derives $\bm y=(g^{\gamma_i})_i$,
 $\bm\nu=(g^{\mu_c})_c$, and the PESAT instance; the bit witness and $O$
 are unchanged.
 \item \textbf{Scope:} the first core F2Z reduction.
\end{itemize}

\item \textbf{Profile restriction: weighted coefficient shape.}
The chunked mod-$q$ profile requires
\[
 u_{c,b,j}=\pi_q(w_b2^j)w'_c,
 \qquad
 0\le w_b<2^{q_{\rm bits}},\qquad w'_c\in\Fq.
\]
This restricts the public coefficient; it does not define another relation.

\end{itemize}

\subsubsection{Active-chunk instantiation of
  \texorpdfstring{$\Pi_1$}{Pi1}}

For the admitted weighted shape, write
\[
 w_b=\sum_{\ell=0}^{L-1}2^{c_w\ell}\omega_b^{(\ell)},
 \qquad
 u_c^{(\ell)}:=\sum_b\omega_b^{(\ell)}\INT(D(b,c)).
\]
The proof fragment is
\begin{equation}
 \pi_1^{\rm chunk}
 :=\bigl(\nu^{(\ell)},u^{(\ell)}\bigr)_{
       \ell\in\mathcal A\uparrow},
 \qquad
 \begin{aligned}
  u^{(\ell)}&=(u_c^{(\ell)})_{c\in\bitsset^s},\\
  \nu^{(\ell)}&=(g^{u_c^{(\ell)}})_{c\in\bitsset^s}.
 \end{aligned}
 \label{eq:pi1-chunk-proof}
\end{equation}
For each active chunk in increasing $\ell$, $\nu^{(\ell)}$ is a NARG record
and $u^{(\ell)}$ is a checked, non-absorbed hint.  Inactive chunks have the
derived value $u^{(\ell)}=0$.  The verifier rejects a missing, duplicate,
out-of-order, or noncanonical active pair; requires
$0\le u_c^{(\ell)}<2^{127}$; and checks
$g^{u_c^{(\ell)}}=\nu_c^{(\ell)}$ before its next squeeze.  Because
$\operatorname{ord}(g)=2^{128}-1$, the stated range makes the accepted exponent
unique.  The checked values reconstruct the paper's column messages:
\begin{equation}
 \mu_c
 :=\sum_{\ell=0}^{L-1}2^{c_w\ell}u_c^{(\ell)}
 =\sum_b w_b\INT(D(b,c)).
 \label{eq:chunk-reconstructs-mu}
\end{equation}
Inside $\Pi_1$, immediately after receiving the final active vector, the
verifier checks
\begin{equation}
 \sum_c w'_c\,\pi_q(\mu_c)=y.
 \label{eq:pi1-weighted-reconstruction}
\end{equation}
This is the tensor-reconstruction check of
Equation~\eqref{eq:pi1-paper-reconstruction}; it creates no named successor
relation and leaves no public scalar obligation for the opening phase.

For each active $\ell$, the verifier also derives
\[
i=(b\ll w)\mid j,\qquad
\omega_i^{(\ell)}=\omega_b^{(\ell)}2^j,\qquad
 y_i^{(\ell)}=g^{\omega_i^{(\ell)}}.
\]
Thus the single paper target is realized by the profile family
\begin{equation}
 \bigwedge_{\ell\in\mathcal A}
 \RPESAT(\K,Q_{\Ftwo\!\Z}^{(\ell)},\mathcal C,\Pack).
 \label{eq:chunk-pesat-family}
\end{equation}

\subsection{The encoded F2Z PESAT relation}

For one active profile chunk, define
\[
 Q_c^{(\ell)}(\Bbit)
 :=\prod_i\left(1+(y_i^{(\ell)}-1)\Bbit[c,i]\right)
   -\nu_c^{(\ell)}.
\]
The corresponding named paper relation is the encoded PESAT instance
\begin{align}
 &\RPESAT(\K,Q_{\Ftwo\!\Z}^{(\ell)},\mathcal C,\Pack)
 \notag\\[-2pt]
 &\quad:=
 \Bigl\{((\bm y^{(\ell)},\bm\nu^{(\ell)}),O;\Bbit):
 \Bbit\in\Ftwo^{2^s\times2^d},\ 
 O=\Enc_{\mathcal C}(\Pack(\Bbit)),\ 
 Q_c^{(\ell)}(\Bbit)=0\ (\forall c)\Bigr\}.
 \label{eq:rel-gp}
\end{align}
Omitting $\ell$ gives the paper's unchunked notation
$\RPESAT(\K,Q_{\Ftwo\!\Z},\mathcal C,\Pack)$.  The profile's ``grand
product'' is this PESAT relation, not a distinct formal relation.

\paragraph{Direct reductions and classifications.}
\begin{itemize}

\item \textbf{Reduction 1: Linearize with GKR.}
\[
 \RPESAT(\K,Q_{\Ftwo\!\Z},\mathcal C,\Pack)
 \xRightarrow{\Pi_{\rm red}}
 \RLin(\K,\mathcal C,\Pack).
\]
For the chunked family, $\Pi_{\rm red}$ is applied once per
$\ell\in\mathcal A$.
\begin{itemize}
 \item \textbf{Input:} the chunk's PESAT instance, the unchanged root
 $\rho=\Root(O)$, and the transcript after $\zeta^{(\ell)}$ and
 $e_0^{(\ell)}$; call this bundle $\mathcal I_{\rm GKR}^{(\ell)}$.
 \item \textbf{Proof:}
 \[
  (\tr',\mathit{pt}_\ell,\mu_\ell)
  \leftarrow
  \mathsf{GKR.Verify}
  (\mathsf{pp}_{\rm GKR},\mathcal I_{\rm GKR}^{(\ell)},
   \pi_{\rm red}^{(\ell)};\tr).
 \]
 $\pi_{\rm red}^{(\ell)}$ is one canonical GKR fragment.  Its profile fixes the
 messages, challenge order, encodings, tags, and rejection rules; the call
 consumes the whole fragment and returns $\tr'$.
 \item \textbf{Result:} the verifier derives $(\bm a_\ell,\mu_\ell)$ for
 $\RLin(\K,\mathcal C,\Pack)$, retaining the same $O$ and bit witness.  In the
 current profile, with
 $B_{\mathrm{bit},\K}:=\iota_{2\to\K}(\Bbit)$ coordinatewise,
 \[
  \bm a_\ell[z]=\eq(z,\mathit{pt}_\ell),
  \qquad
  \MLE{B_{\mathrm{bit},\K}}(\mathit{pt}_\ell)=\mu_\ell.
 \]
 This equation defines the LIN successor; it is not another paper relation.
 \item \textbf{Scope:} this edge is solid once the GKR profile fixes a
 byte-exact transcript and proves the reduction.  Its error is
 $\varepsilon_{\rm red}^{(\ell)}$ per active chunk.
\end{itemize}

\item \textbf{Classification: Degree-zero, zero-ideal PIMSAT.}
With $S=\K[X]$, $S_{\le0}=\K$, and $I=(0)$, membership means exactly
$Q_c^{(\ell)}=0$.  Boolean witnesses keep their bit domain; unrestricted
$\K$ variables also require $Z^2-Z=0$.

\end{itemize}

\subsection{The encoded LIN relation}

Let $\iota_{2\to\K}:\Ftwo\hookrightarrow\K$ act coordinatewise and put
$\bm b_{\K}:=\iota_{2\to\K}(\pi_2(\bm b))$.  The paper's linear target is
\begin{align}
 \RLin(\K,\mathcal C,\Pack):=
 \Bigl\{((\bm a,\mu),O;\bm b):\;&
 \bm b\in\bitsset^n,
 \quad\bm a\in\K^n,
 \quad\mu\in\K,\notag\\[-2pt]
 &O=\Enc_{\mathcal C}(\Pack(\pi_2(\bm b))),
 \quad\langle\bm a,\bm b_{\K}\rangle_{\K}=\mu\Bigr\}.
 \label{eq:rel-lin-pack}
\end{align}

\paragraph{Direct reductions.}
\begin{itemize}

\item \textbf{Reduction 1: Run the linear-relation IOP.}
\[
 \RLin(\K,\mathcal C,\Pack)
 \xRightarrow{\Pi_{\rm iopp}}
 \mathcal R_{\rm triv}.
\]
\begin{itemize}
 \item \textbf{Proof:} $\pi_{\rm iopp}$ proves the derived $(\bm a,\mu)$ claim
 against the public root $\rho=\Root(O)$.
 \item \textbf{Result:} acceptance or rejection; no public relation remains.
 \item \textbf{Scope:} the final solid paper-level edge.  The profile below
 gives one concrete implementation.
\end{itemize}

\end{itemize}

\subsubsection{Current profile's internal expansion of
  \texorpdfstring{$\Pi_{\rm iopp}$}{Pi-iopp}}

The chunked recursive-opening profile implements $\Pi_{\rm iopp}$ on the
MLE-shaped LIN family emitted by $\Pi_{\rm red}$ as follows.  These are
internal claim transformations, not additional solid edges in the paper's
relation graph.
\begin{enumerate}

\item \textbf{Per-claim ring switch.}
For each active $\ell$, the prover transmits
\[
 s^{(\ell)}=(s_v^{(\ell)})_{v<128},
 \qquad
 s_v^{(\ell)}:=\MLE{B_{\mathrm{bit},v}}(r_{\hi}^{(\ell)}).
\]
After all such vectors are absorbed, the verifier squeezes one shared $r''$
and derives a packed claim
\begin{equation}
 O=\Enc_{\mathcal C}(P),
 \qquad
 \langle B_\ell,P\rangle_{\K}=\beta_\ell.
 \label{eq:profile-packed-claim}
\end{equation}
The witness is $P=\Pack(\pi_2(\bm b))$ and $\rho=\Root(O)$ is unchanged.  This step
has error $\varepsilon_{\rm ring}^{(\ell)}$ on the admitted MLE-shaped source;
it does not specify a general ring switch for arbitrary LIN coefficients.

\item \textbf{Batch.}
The verifier squeezes fresh
$\bm\eta\in\K^{|\mathcal A|}$ and derives, without a proof field,
\[
 B:=\sum_{\ell\in\mathcal A}\eta_\ell B_\ell,
 \qquad
 \beta:=\sum_{\ell\in\mathcal A}\eta_\ell\beta_\ell.
\]
The resulting internal claim is
$O=\Enc_{\mathcal C}(P)$ and
$\langle B,P\rangle_{\K}=\beta$, with batching error
$\varepsilon_{\rm batch}$.

\item \textbf{Recursive opening.}
The prover transmits $\pi_{\rm open}=\pi_{\rm lig}$, one profile-defined
recursive constrained-code proof containing its internal proximity messages
and authenticated openings.  The verifier derives all opening challenges and
accepts only if the proof verifies against $(\rho,B,\beta)$.  This completes
$\Pi_{\rm iopp}$ with error $\varepsilon_{\rm open}$.

\end{enumerate}

The intended paper ring-switch construction is not used as an independent
solid edge here: the three operations above are the selected profile's
internal implementation of the already named
$\RLin(\K,\mathcal C,\Pack)\xRightarrow{\Pi_{\rm iopp}}\mathcal R_{\rm triv}$
edge.

\subsection{The virtual Lin-BitsDual extension}

For a public $\Ftwo$-linear map $T:\Ftwo^n\to\Ftwo^{n'}$ and
$h_0\in\Ftwo^{n'}$, let
$\operatorname{lift}_2:\Ftwo\to\{0,1\}\subset\Z$ act coordinatewise and put
$\bar h=h_0+T(\pi_2(\bm b))$.  Define
\begin{align}
 \RLBD^{\virt}:=\Bigl\{((T,\bm a_q,\mu,h_0),O;\bm b):\;&
 \bm b\in\bitsset^n,
 \quad\bm a_q\in\Fq^{n'},
 \quad\mu\in\Fq,\notag\\[-2pt]
 &O=\Enc_{\mathcal C}(\Pack(\pi_2(\bm b))),
 \quad
 \langle\pi_q(\operatorname{lift}_2(\bar h)),\bm a_q\rangle_{\Fq}
 =\mu\Bigr\}.
 \label{eq:rel-lbd-virtual}
\end{align}

\paragraph{Direct reductions.}
\begin{itemize}

\item \textbf{Conditional virtual reduction.}
\[
 \RLBD^{\virt}
 \overset{\Pi_{\rm partial\text{-}core}}{\dashrightarrow}
 \mathcal R_{h\text{-}{\rm Lin}}^{\virt}.
\]
\begin{itemize}
 \item \textbf{Proof:} $\pi_{\rm virtual}$ carries the chosen partial-core and
 GKR messages; the virtual table is not transmitted.
 \item \textbf{Result:} the verifier derives $(\bm a_{\K},\mu')$.
 \item \textbf{Scope:} available only after a virtual profile defines the
 proof format, transcript, verifier, and error bound.
\end{itemize}

\end{itemize}

Let $h_{0,\K}:=\iota_{2\to\K}(h_0)$ and let $T_{\K}$ be the scalar extension
of $T$.  Only $\bm b$ is committed.  The virtual linear seam is
\begin{align}
 \mathcal R_{h\text{-}{\rm Lin}}^{\virt}:=
 \Bigl\{((T,h_0,\bm a_{\K},\mu'),O;\bm b):\;&
 \bm b\in\bitsset^n,
 \quad\bm a_{\K}\in\K^{n'},
 \quad\mu'\in\K,\notag\\[-2pt]
 &O=\Enc_{\mathcal C}(\Pack(\pi_2(\bm b))),\notag\\[-2pt]
 &\langle\bm a_{\K},h_{0,\K}+T_{\K}(\bm b_{\K})\rangle_{\K}
 =\mu'\Bigr\}.
 \label{eq:rel-virtual-lin}
\end{align}
Its public adjoint rewrite transmits nothing:
\[
 \mathcal R_{h\text{-}{\rm Lin}}^{\virt}
 \equiv_{\mathsf{Adj}_{T_{\K}^{*}}}
 \RLin(\K,\mathcal C,\Pack),
\]
\[
 \bm a_b:=T_{\K}^{*}(\bm a_{\K}),
 \qquad
 \mu_b:=\mu'-\langle\bm a_{\K},h_{0,\K}\rangle_{\K}.
\]
This exact rewrite need not preserve the MLE coefficient shape admitted by the
current internal ring switch.  A virtual profile must therefore supply a
compatible $\Pi_{\rm iopp}$ adapter or opening; the direct profile rejects the
virtual source.

\subsection{The complete relation graph and proof object}

Let $\mathcal R_{\rm triv}:=\{(\bot,\bot;\bot)\}$.  The paper-level relation
graph is
\begin{equation}
\begin{aligned}
 \RPIMSAT^{\rm app}
 &\xRightarrow{\mathsf{Encode}}
 \RPIMSAT^{\rm enc},
 \\[1pt]
 \RPIMSAT^{\rm enc}
 &\xRightarrow{\Pi_{{\rm app},1}}
 \RPESAT^{\rm app}
 \xRightarrow{\Pi_{{\rm app},2}}
 \RLMFB,
 \\[1pt]
 \RLMFB
 &\equiv_{\mathsf{Bitify}}
 \RLBD
 \xRightarrow{\Pi_1}
 \RPESAT(\K,Q_{\Ftwo\!\Z},\mathcal C,\Pack)
 \\[1pt]
 \RPESAT(\K,Q_{\Ftwo\!\Z},\mathcal C,\Pack)
 &\xRightarrow{\Pi_{\rm red}}
 \RLin(\K,\mathcal C,\Pack)
 \xRightarrow{\Pi_{\rm iopp}}
 \mathcal R_{\rm triv}.
\end{aligned}
\label{eq:solid-reduction-graph}
\end{equation}
$\mathsf{Encode}$ and Bitify add no proof fields; chunking and recursive opening
refine the displayed edges.  The semantic proof object is
\begin{equation}
 \pi_{\mathrm{F2Z}}
 =\left(
   \pi_{\rm PIOP},
   \bigl(\nu^{(\ell)},u^{(\ell)},\pi_{\rm red}^{(\ell)}\bigr)_{
         \ell\in\mathcal A\text{ increasing}},
   (s^{(\ell)})_{\ell\in\mathcal A\text{ increasing}},
   \pi_{\rm lig}
 \right).
 \label{eq:proof-object}
\end{equation}
Its replay schedule is specified in Section~\ref{sec:interaction}.  Public
inputs are external to the proof; other omitted values are derived.  The v1
byte profile partitions transmitted fields, in replay order, as
\begin{equation}
 \pi_{\rm wire}:=(N_{\rm F2Z},H_{\rm F2Z}),\qquad
 \begin{aligned}
  \NargLen(\pi)&:=|N_{\rm F2Z}|,&
  \HintLen(\pi)&:=|H_{\rm F2Z}|,\\
  \PayloadLen(\pi)&:=|N_{\rm F2Z}|+|H_{\rm F2Z}|,&
    \WireLen(\pi)&:=
    |\HostEncode(N_{\rm F2Z},H_{\rm F2Z})|.
 \end{aligned}
 \label{eq:proof-channel-lengths}
\end{equation}
Benchmarks \must{} report all four lengths.  Section~\ref{sec:wire-container}
fixes their byte meaning, the outer 72-byte container, and the exact record
counts and exhaustion rules.

With edge errors $\varepsilon_{{\rm app},1}$, $\varepsilon_{{\rm app},2}$,
$\varepsilon_{\rm red}^{(\ell)}$, $\varepsilon_{\rm ring}^{(\ell)}$,
$\varepsilon_{\rm batch}$, and $\varepsilon_{\rm open}$,
\[
 \varepsilon_{\rm F2Z}
 \le \varepsilon_{{\rm app},1}+\varepsilon_{{\rm app},2}
    +\sum_{\ell\in\mathcal A}
       \bigl(\varepsilon_{\rm red}^{(\ell)}
             +\varepsilon_{\rm ring}^{(\ell)}\bigr)
    +\varepsilon_{\rm batch}+\varepsilon_{\rm open}.
\]

\endgroup

\section{Prover--verifier interaction}
\label{sec:interaction}

\definecolor{pvnavy}{HTML}{17365D}
\definecolor{pvtraffic}{HTML}{F3F4F6}
\definecolor{pvhint}{HTML}{A15C00}

\newcommand{\Nbadge}[1]{\textcolor{specblue}{\chN[#1]}}
\newcommand{\Hbadge}[1]{\textcolor{pvhint}{\chH[#1]}}
\newcommand{\Dbadge}[1]{\textcolor{specmuted}{\chD[#1]}}
\newcommand{\Splitbadge}[1]{%
  \textcolor{specblue}{\chN}\!/
  \textcolor{pvhint}{\chH}[#1]}

\newcommand{\idCS}{$\mathsf{CS}[h]$}
\newcommand{\idLINH}{$\mathsf{LIN}[h]$}
\newcommand{\idSUMIN}{$\mathsf{SC}_{\rm in}[H]$}
\newcommand{\idPESAT}{$\mathsf{PESAT}$}
\newcommand{\idSUMH}{$\mathsf{SUM}[H]$}
\newcommand{\idEVALH}{$\mathsf{EVAL}[H]$}
\newcommand{\idLINPACK}{$\mathsf{LIN}[f]$}
\newcommand{\idLINID}{$\mathsf{LIN}[\mathsf{id}]$}

\newcolumntype{L}{>{\raggedright\arraybackslash}p{2.04in}}
\newcolumntype{T}{>{\centering\arraybackslash\columncolor{pvtraffic}}p{0.65in}}
\newcolumntype{R}{>{\raggedright\arraybackslash}p{2.24in}}

\newenvironment{pvexchange}
  {\par\scriptsize
   \setlength{\tabcolsep}{3pt}%
   \setlength{\abovedisplayskip}{2pt}%
   \setlength{\belowdisplayskip}{2pt}%
   \renewcommand{\arraystretch}{1.02}%
   \begin{tabular}{@{}LTR@{}}
   \rowcolor{pvnavy}
   \multicolumn{1}{c}{\color{white}\bfseries Prover $\mathsf P$}
   &\multicolumn{1}{c}{\color{white}\bfseries Traffic}
   &\multicolumn{1}{c}{\color{white}\bfseries Verifier $\mathsf V$}
   \\
   \addlinespace[1pt]}
  {\end{tabular}\par}

\newenvironment{pvlongexchange}
  {\par\scriptsize
   \setlength{\LTleft}{0pt}\setlength{\LTright}{\fill}%
   \setlength{\LTpre}{0pt}\setlength{\LTpost}{0pt}%
   \setlength{\tabcolsep}{3pt}%
   \setlength{\abovedisplayskip}{2pt}%
   \setlength{\belowdisplayskip}{2pt}%
   \renewcommand{\arraystretch}{1.02}%
   \begin{longtable}{@{}LTR@{}}
   \rowcolor{pvnavy}
   \multicolumn{1}{c}{\color{white}\bfseries Prover $\mathsf P$}
   &\multicolumn{1}{c}{\color{white}\bfseries Traffic}
   &\multicolumn{1}{c}{\color{white}\bfseries Verifier $\mathsf V$}
   \\
   \addlinespace[1pt]
   \endfirsthead
   \rowcolor{pvnavy}
   \multicolumn{1}{c}{\color{white}\bfseries Prover $\mathsf P$}
   &\multicolumn{1}{c}{\color{white}\bfseries Traffic}
   &\multicolumn{1}{c}{\color{white}\bfseries Verifier $\mathsf V$}
   \\
   \addlinespace[1pt]
   \endhead}
  {\end{longtable}\par}

\newcommand{\Pmsg}[1]{\(\xrightarrow{\;\scriptstyle #1\;}\)}
\newcommand{\Vmsg}[1]{\(\xleftarrow{\;\scriptstyle #1\;}\)}
\newcommand{\Pturn}[1]{#1 & & \\[1pt]}
\newcommand{\Vturn}[1]{& & #1 \\[1pt]}
\newcommand{\Psend}[1]{& \Pmsg{#1} & \\[-1pt]}
\newcommand{\Vsend}[1]{& \Vmsg{#1} & \\[-1pt]}

\newenvironment{pvstep}[3][2.0in]
  {\Needspace{#1}%
   \par\addvspace{0.45\baselineskip}%
   \refstepcounter{clgstep}%
   \clgmarkstart
   \noindent
   {\normalsize\bfseries\sffamily\textcolor{clgtitlec}%
     {P\theclgstep/V\theclgstep.}~#2}%
   \nopagebreak\par\vspace{0.08\baselineskip}%
   #3%
   \clgboard%
   \nopagebreak}
  {\clgreserve\par\addvspace{0.12\baselineskip}}

\subsection{End-to-end proof decomposition and transcript channels}
\label{sec:transcript-channels}

The end-to-end proof consists of the configured PIOP proof followed by the
F2Z proof:
\begin{equation}
 \pi_{\rm e2e}:=\left(
   \pi_{\rm PIOP},
   \pi_{\rm F2Z}
 \right).
 \label{eq:end-to-end-proof-object}
\end{equation}
Verification of $\pi_{\rm PIOP}$ reduces the initial constraint statement to
the F2Z handoff
\begin{equation}
 \left\langle\bm v,\operatorname{vec}(H_q)\right\rangle_{\Fq}=\mu.
 \label{eq:piop-f2z-handoff}
\end{equation}
The pair $(\bm v,\mu)$ is produced by the PIOP verifier and is not repeated as
an F2Z proof field.

Within F2Z, write
\begin{equation}
 \pi_1:=\left(
   \pi_{1,{\rm in}},
   \bm\nu,\bm\mu
 \right),
 \qquad
 \pi_{1,{\rm in}}:=(g_i^{\rm shape})_{i=0}^{m_h-1}.
 \label{eq:pi1-proof-object}
\end{equation}
The semantic F2Z proof object is
\begin{equation}
 \pi_{\mathrm{F2Z}}
 =\left(
   \pi_1,
   \pi_{\rm GKR},\pi_{\rm MLE},
   \bm s,
   \pi_{\rm open}
 \right).
 \label{eq:proof-object}
\end{equation}
The fragment $\pi_{1,{\rm in}}$ contains the round polynomials of the standard
input Sumcheck that reduces Equation~\eqref{eq:piop-f2z-handoff} to
\begin{equation}
 \kappa\widetilde H_q(\bm s_h)=\tau.
 \label{eq:input-sumcheck-output}
\end{equation}
The point $\bm s_h$, scalars $\kappa$ and $\tau$, and factors
$(\bm a,\bm a')$ are verifier-derived and are not additional proof fields.
The ring-switch component is
$\bm s=(\bm s_\theta)_{\theta\in\Theta}$, with each
$\bm s_\theta\in\K^{128}$; the ordered set $\Theta$ and its coefficient
decomposition are public and are not proof fields.

Equations~\eqref{eq:end-to-end-proof-object}--\eqref{eq:proof-object} describe
semantic proof components.  The developer configuration fixes the encoding of
$\pi_{\rm PIOP}$, while Section~\ref{sec:wire-format} fixes the encoding of
$\pi_{\rm F2Z}$.  The F2Z v1 container partitions its transmitted records into
the NARG stream $N_{\rm F2Z}$ and the hint stream $H_{\rm F2Z}$:
\begin{equation}
 \pi_{\rm wire}:=(N_{\rm F2Z},H_{\rm F2Z}),\qquad
 \begin{aligned}
  \NargLen(\pi_{\rm F2Z})&:=|N_{\rm F2Z}|,&
  \HintLen(\pi_{\rm F2Z})&:=|H_{\rm F2Z}|,\\
  \PayloadLen(\pi_{\rm F2Z})&:=|N_{\rm F2Z}|+|H_{\rm F2Z}|,&
  \WireLen(\pi_{\rm F2Z})&:=|\HostEncode(N_{\rm F2Z},H_{\rm F2Z})|.
 \end{aligned}
 \label{eq:proof-channel-lengths}
\end{equation}

Equation~\eqref{eq:proof-channel-lengths} defines the two F2Z proof streams;
the following badges identify how each transmitted value is handled.

\begingroup\footnotesize
\renewcommand{\arraystretch}{1.05}
\begin{tabularx}{\linewidth}{@{}p{0.48in}p{0.90in}X@{}}
\toprule
\textbf{Badge}&\textbf{Source}&\textbf{Verifier action}\\
\midrule
$\Nbadge{x}$&NARG stream&Receive $x$ and verify it, and absorb it into the
verifier's transcript.\\
$\Hbadge{x}$&Hint stream&Receive $x$ and verify it, but do not absorb it into the
verifier's transcript.\\
$\Dbadge{x}$&No proof bytes&Receive nothing; $x$ is public or derived by the
verifier.\\
$\Splitbadge{\pi}$&Both streams&Receive each field from its specified stream.
Absorb NARG fields and verify Hint fields.\\
\bottomrule
\end{tabularx}
\endgroup

The verifier completes every required check and absorption before deriving the
next challenge.

\paragraph{Transcript absorption.}
Every absorbed value uses the encoding
\begin{equation}
 \Absorb\!\left(
   \mathsf{DOMAIN}
   \parallel\mathsf{PAYLOAD\_SIZE}
   \parallel\mathsf{PAYLOAD}
 \right).
 \label{eq:transcript-absorb}
\end{equation}
Section~\ref{sec:wire-format} specifies the canonical encoding of each domain
and payload, together with the transcript replay rules.  The domain identifies
the semantic meaning of the absorbed value.

Right ledgers contain post-step claims; \emph{internal} claims are shown only
when they explain why a later acceptance check cannot yet discharge the
original relation.

\setlength{\textwidth}{5.63in}
\setlength{\columnwidth}{5.63in}
\setlength{\linewidth}{5.63in}
\setlength{\hsize}{5.63in}
\setlength{\marginparwidth}{1.72in}
\setlength{\marginparsep}{0.10in}

% --------------------------------------------------------------------------
\begin{pvstep}{Bind the constraint statement and commitment}{%
  \clgnewrel{\idCS}
    {$(x;H_q)\in\mathcal R_{\rm CS}$}%
  \clgdef{$\rho$}{$\rho=\Root(O)$}%
}
\begin{pvexchange}
\Pturn{Construct the committed and PIOP-facing witness views:
\(\begin{aligned}
 \pi_2(\bm h)&=\bm M\pi_2(\bm f),\\
 \operatorname{vec}(H)&=\pi_2(\bm h),\\
 \operatorname{vec}(F)&=\pi_2(\bm f),\qquad P:=\Pack(F),\\
 O&:=\Enc_{\mathcal C}(P),\qquad \rho:=\Root(O).
\end{aligned}\)
and retain $(\bm h,H,F,P,O)$.}
\Psend{\Dbadge{x,\rho}}
\Vturn{Validate the constraint statement, the developer configuration, and
the commitment root.  Bind $(x,\rho)$ as required by the configured PIOP.}
\end{pvexchange}
\end{pvstep}

\smallskip
\noindent\textbf{Configured matrix indexing.}
The developer configuration fixes a bijection between each Boolean address
$z\in\bitsset^{m_h}$ and a pair
\[
 z=(r,c)\in\bitsset^{d_r}\times\bitsset^{d_c},
 \qquad m_h=d_r+d_c.
\]
Using this fixed address split, define
\[
 H[r,c]:=\bigl[\pi_2(\bm h)\bigr]_z,
 \qquad
 \operatorname{vec}(H)=\pi_2(\bm h).
\]
Thus $H$ is a matrix view of the same witness bits, not a new witness.  The
tables $H_{\mathbb Z}$, $H_q$, and $H_{\K}$ contain these same entries embedded
into $\mathbb Z$, $\Fq$, and $\K$, respectively.  We use brackets for table
entries, such as $H_q[r,c]$, and parentheses only for multilinear-extension
evaluations, such as $\widetilde H_q(\bm x_r,\bm x_c)$.

% --------------------------------------------------------------------------
\begin{pvstep}[2.65in]{Run the PIOP}{%
  \clgedge[PIOP]{$\Pi_{\rm PIOP}$}
    {$\mathcal R_{\rm CS}$}
    {$\mathcal R_{\rm LinBitsDual}[\bm h]$}%
  \clgdone{\idCS}%
  \clgnewrel{\idLINH}
    {$\langle\bm v,\operatorname{vec}(H_q)\rangle_{\Fq}=\mu$}%
}
\begin{pvexchange}
\Pturn{Run the configured PIOP on the constraint statement and witness:
\[
 \pi_{\rm PIOP}\leftarrow
 \Pi_{\rm PIOP}.\mathsf{Prove}
   (\mathsf{pp}_{\rm PIOP},x,\operatorname{vec}(H_q)).
\]}
\Psend{\pi_{\rm PIOP}}
\Vturn{Run the matching PIOP verifier:
\[
 (\mathsf{ok}_{\rm PIOP},\bm v,\mu)\leftarrow
 \Pi_{\rm PIOP}.\mathsf{Verify}
   (\mathsf{pp}_{\rm PIOP},x,\pi_{\rm PIOP}).
\]
Require $\mathsf{ok}_{\rm PIOP}=1$.  Its output is the F2Z handoff
\[
 \boxed{\left\langle\bm v,\operatorname{vec}(H_q)\right\rangle_{\Fq}=\mu}.
\]}
\Vsend{\Dbadge{\bm v,\mu}}
\end{pvexchange}
\end{pvstep}

% --------------------------------------------------------------------------
\begin{pvstep}[3.05in]{Reduce the PIOP handoff with Sumcheck}{%
  \clgedge[Sumcheck]{$\Pi_{1,{\rm in}}$}
    {$\mathcal R_{\rm LinBitsDual}[\bm h]$}
    {$\kappa\widetilde H_q(\bm s_h)=\tau$}%
  \clgdone{\idLINH}%
  \clgnewinternal{\idSUMIN}
    {$\kappa\widetilde H_q(\bm s_h)=\tau$}%
}
\begin{pvexchange}
\Vturn{Initialize the F2Z transcript by binding its domain, handoff statement,
configuration, and commitment root.}
\Pturn{Run the standard $m_h$-round Sumcheck on the Lin-BitsDual claim from
Section~\ref{sec:linbits-instance}:
\[
 \begin{aligned}
  \bm X&=(\bm X_r,\bm X_c),\\
  G_{\rm shape}(\bm X)
    &:=\widetilde{\bm v}(\bm X_r,\bm X_c)\\[-2pt]
    &\phantom{:=}{}\cdot\widetilde H_q(\bm X_r,\bm X_c),\\
  \mu&=\sum_{r,c}G_{\rm shape}(r,c),\\
  \pi_{1,{\rm in}}&:=(g_i^{\rm shape})_{i=0}^{m_h-1}.
 \end{aligned}
\]}
\Psend{\Nbadge{\pi_{1,{\rm in}}}}
\Vturn{Verify the Sumcheck transcript.  It yields a terminal point
$\bm s_h=(\bm s_r,\bm s_c)$ and value $\tau$ satisfying
\[
 \tau=\widetilde{\bm v}(\bm s_r,\bm s_c)
      \widetilde H_q(\bm s_r,\bm s_c),
 \qquad
 \kappa:=\widetilde{\bm v}(\bm s_r,\bm s_c),
\]
and derive
\(\begin{aligned}
 a_r&:=\eq_{\Fq}(r,\bm s_r),\\
 a'_c&:=\kappa\eq_{\Fq}(c,\bm s_c).
\end{aligned}\)
The resulting terminal claim is
\[
 \boxed{
  \kappa\widetilde H_q(\bm s_r,\bm s_c)
  =\tau
  =\sum_c a'_c\sum_r a_rH_q[r,c]
 }.
\]
Absorb the derived seam and derive the canonical integer lifts
$\gamma_r:=\operatorname{lift}_q(a_r)$.}
\Vsend{\Dbadge{\bm s_h,\kappa,\tau}}
\end{pvexchange}
\end{pvstep}

% --------------------------------------------------------------------------
\begin{pvstep}[3.15in]{Fold each column and encode the folds as product claims}{%
  \clgdone{\idSUMIN}%
  \clgnewrel{\idPESAT}{$\mathcal R_{\rm PESAT}[\bm h]$}%
  \clgnote{\idPESAT}{%
    $\nu_c=\prod_r
      \bigl(1+(g^{\gamma_r}-1)H_{\K}[r,c]\bigr)$}%
  \clgdef{$P$}{$P=\Pack(F)$}%
}
\begin{pvlongexchange}
\Pturn{For every column $c$, compute
\(\begin{aligned}
 \widehat\mu_c&:=\sum_r\gamma_rH_{\mathbb Z}[r,c],\\
 \mu_c&:=\widehat\mu_c,\\
 \nu_c&:=g^{\mu_c}\in\K^\times.
\end{aligned}\)}
\Psend{\Nbadge{\bm\nu}}
\Psend{\Hbadge{\bm\mu}}
\Vturn{Before deriving any dependent challenge, require
\[
 |\bm\mu|=2^{d_c},
\]
and, for every $c$,
\[
 0\le\mu_c<2^{127},\qquad g^{\mu_c}=\nu_c.
\]
Then close the retained Sumcheck claim by checking
\[
 \boxed{\sum_c a'_c\pi_q(\mu_c)=\tau}.
\]}
\Vturn{After accepting the hint, derive
\(\begin{aligned}
 \zeta&\leftarrow\Squeeze_{\K}^{d_c}(\tr),\\
 e_0&:=\MLE{\bm\nu}(\zeta).
\end{aligned}\)}
\Vsend{\Dbadge{\zeta}}
\Vturn{Record the successor claim for P5/V5:
\(\begin{aligned}
 Q_{r,c}&:=1+(g^{\gamma_r}-1)H_{\K}[r,c],\\
 \nu_c&=\prod_r Q_{r,c}.
\end{aligned}\)
This equality is not checked in P4/V4, and no entry of $H_{\K}$ is sent.}
\end{pvlongexchange}
\end{pvstep}

% --------------------------------------------------------------------------
\begin{pvstep}[2.35in]{Run the GKR reduction}{%
  \clgedge[GKR subprotocol]{$\Pi_{\rm GKR}$}
    {$\mathcal R_{\rm PESAT}[\bm h]$}
    {$\mathcal R_{\rm sum}^{h}$}%
  \clgdone{\idPESAT}%
  \clgnewrel{\idSUMH}
    {$S(H_{\K})=\sigma$}%
  \clgnote{\idSUMH}{%
    $S(H_{\K})=\sum_z\omega_{\xi,z}H_{\K}[z]$}%
}
\begin{pvexchange}
\Pturn{Set
\(\displaystyle
 \mathcal I_{\rm GKR}:=
 \mathsf{GKRInput}(\bm\gamma,\bm\nu,\zeta,e_0)
\)
and compute
\(\displaystyle
 \pi_{\rm GKR}\leftarrow
 \mathsf{GKR.Prove}
 (\mathsf{pp}_{\rm GKR},\mathcal I_{\rm GKR},H_{\K};\tr)\).}
\Psend{\Splitbadge{\pi_{\rm GKR}}}
\Vturn{Compute
\(\begin{aligned}
 (\tr',\bm\omega_\xi,\sigma)&\leftarrow
   \mathsf{GKR.Verify}\\
 &\quad(\mathsf{pp}_{\rm GKR},\mathcal I_{\rm GKR},
        \pi_{\rm GKR};\tr),
\end{aligned}\)
require success and the exact subordinate endpoint, then set
$\tr\leftarrow\tr'$.  The successor is
\(\displaystyle
 \sum_{z\in\bitsset^{m_h}}
   \omega_{\xi,z}H_{\K}[z]=\sigma.
\)}
\end{pvexchange}
GKR operates on $H_{\K}$ and does not query $O$.
\end{pvstep}

% --------------------------------------------------------------------------
\begin{pvstep}[3.05in]{Reduce the GKR sum claim to an MLE claim}{%
  \clgedge[row LinCheck]{$\Pi_{\rm MLE}$}
    {$\mathcal R_{\rm sum}^{h}$}
    {$\mathcal R_{\rm eval}^{h}$}%
  \clgdone{\idSUMH}%
  \clgnewrel{\idEVALH}
    {$\widetilde H_{\K}(\bm r_h)=\eta$}%
  \clgundef{$P$}%
}
\begin{pvexchange}
\Pturn{Run the standard $m_h$-round Sumcheck on the weighted-sum claim:
\[
 \begin{aligned}
  G_{\rm MLE}(\bm X)
    &:=\widetilde{\bm\omega_\xi}(\bm X)
       \widetilde H_{\K}(\bm X),\\
  \sigma&=\sum_{z\in\bitsset^{m_h}}G_{\rm MLE}(z),\\
  \pi_{\rm MLE}
    &:=\bigl((g_i^{\rm MLE})_{i=0}^{m_h-1},\eta\bigr),
 \end{aligned}
\]
where the prover supplies
$\eta:=\widetilde H_{\K}(\bm r_h)$ at the transcript-derived terminal point
$\bm r_h$.}
\Psend{\Splitbadge{\pi_{\rm MLE}}}
\Vturn{Verify the Sumcheck transcript.  It yields a terminal point
$\bm r_h$ and value $\theta$.  Compute
\[
 \lambda:=\widetilde{\bm\omega_\xi}(\bm r_h)
\]
and require the terminal check
\[
 \boxed{\theta=\lambda\eta}.
\]
The resulting successor claim is
\[
 \boxed{\widetilde H_{\K}(\bm r_h)=\eta}.
\]}
\end{pvexchange}
The row LinCheck operates on $H_{\K}$ and does not query $O$.
\end{pvstep}

% --------------------------------------------------------------------------
\begin{pvstep}[2.25in]{Rewrite $\langle\bm a,\bm M\bm f\rangle$ as
$\langle\bm M^{\mathsf T}\bm a,\bm f\rangle$}{%
  \clgedge[adjoint identity]
    {$\begin{gathered}
       \langle\bm a,\bm M_{\K}\bm f_{\K}\rangle_{\K}\\[-2pt]
       =\langle\bm M_{\K}^{\mathsf T}\bm a,\bm f_{\K}\rangle_{\K}
     \end{gathered}$}
    {$\mathcal R_{\rm eval}^{h}$}
    {$\LIN(\K,\mathcal C,\Pack)$}%
  \clgdone{\idEVALH}%
  \clgnewrel{\idLINPACK}{$\langle\bm a_f,\bm f_{\K}\rangle_{\K}=\eta$}%
}
\begin{pvlongexchange}
\multicolumn{3}{@{}p{5.05in}@{}}{
Both parties set $\bm a[z]:=\eq(z,\bm r_h)$ and use
$H_{\K}=\bm M_{\K}\bm f_{\K}$ to rewrite
\[
 \eta
 =\langle\bm a,H_{\K}\rangle_{\K}
 =\langle\bm a,\bm M_{\K}\bm f_{\K}\rangle_{\K}
 =\langle\bm M_{\K}^{\mathsf T}\bm a,\bm f_{\K}\rangle_{\K}.
\]
Set $\bm a_f:=\bm M_{\K}^{\mathsf T}\bm a$.  No message or challenge is
needed.  In the configured packing order, both parties also view this same
coefficient as
\[
 A_f[v,j]:=\bm a_f[128j+v].
\]
This is only a change of coordinates; the claim remains
$\langle\bm a_f,\bm f_{\K}\rangle_{\K}=\eta$.}
\end{pvlongexchange}
\end{pvstep}

% --------------------------------------------------------------------------
\begin{pvstep}[4.35in]{Ring-switch the adjoint claim into the committed packing}{%
  \clgedge[$\Pi_{\rm iopp}$ internal]{$\Pi_{\rm ring}$}
    {$\LIN(\K,\mathcal C,\Pack)$}
    {$\LIN(\K,\mathcal C,\mathsf{id})$}%
  \clgdone{\idLINPACK}%
  \clgnewrel{\idLINID}{$\langle B,P\rangle_{\K}=\beta$}%
  \clgdef{$P$}{$P=\Pack(F)$}%
}
\begin{pvlongexchange}
\multicolumn{3}{@{}p{5.05in}@{}}{
From the public coefficient $A_f$, both parties derive the canonical ordered
decomposition
\[
 A_f[v,j]=\sum_{\theta\in\Theta}A_\theta[v]B_\theta[j]
\]
from Section~\ref{sec:ring-switch-coefficients}.  This represents the same
inner product produced by P7/V7; it is not another reduction or proof claim.}
\\[2pt]
\Pturn{For every $\theta\in\Theta$ and $0\le v<128$, compute
\[
 s_{\theta,v}:=\sum_{j<2^{m_p}}B_\theta[j]F_v[j].
\]}
\Psend{\Nbadge{\bm s_\theta}}
\Vturn{Require $|\bm s_\theta|=128$ for every $\theta$ and anchor all
messages to the adjoint claim:
\[
 \boxed{
  \sum_{\theta\in\Theta}\sum_{v=0}^{127}
    A_\theta[v]s_{\theta,v}=\eta
 }.
\]}
\Vturn{After accepting every $\bm s_\theta$, derive
\(\displaystyle
 r''\leftarrow\Squeeze_{\K}^{7}(\tr).
\)}
\Vsend{\Dbadge{r''}}
\multicolumn{3}{@{}>{\raggedright\arraybackslash}p{5.05in}@{}}{
With $[x]_u$ denoting coordinate $u$ in the fixed $\K/\Ftwo$ basis, both
parties derive, for every $\theta$,
\(\begin{aligned}
 t_{\theta,u}&:=\sum_v[s_{\theta,v}]_u\beta_v,\\
 \Phi_{r''}(x)&:=\sum_u\eq(u,r'')[x]_u,\\
 \beta_\theta&:=\sum_u\eq(u,r'')t_{\theta,u},\\
 \omega_\theta[j]&:=\Phi_{r''}(B_\theta[j]),
\end{aligned}\)
which gives $\langle\bm\omega_\theta,P\rangle_{\K}=\beta_\theta$.}
\\[2pt]
\Vturn{For $\theta\in\Theta$ in order, derive
\(\displaystyle
 \lambda_\theta\leftarrow\Squeeze_{\K}(\tr).
\)}
\Vsend{\Dbadge{(\lambda_\theta)_{\theta\in\Theta}}}
\Pturn{Set
\(\displaystyle
 B:=\sum_{\theta\in\Theta}\lambda_\theta\bm\omega_\theta,
 \qquad
 \beta:=\sum_{\theta\in\Theta}\lambda_\theta\beta_\theta.
\)}
\Vturn{Independently derive the same $(B,\beta)$.}
\end{pvlongexchange}
\end{pvstep}

% --------------------------------------------------------------------------
\begin{pvstep}[2.65in]{Open the packed claim}{%
  \clgedge[$\Pi_{\rm iopp}$ internal]{$\Pi_{\rm open}$}
    {$\LIN(\K,\mathcal C,\mathsf{id})$}
    {$\mathcal R_{\rm triv}$}%
  \clgdone{\idLINID}%
  \clgundef{$P$}%
  \clgundef{$\rho$}%
}
\begin{pvlongexchange}
\Pturn{Absorb the opening target $(m_p,\beta)$.}
\Vturn{Retain the locally derived $B$ from P8/V8 and absorb the same opening
target.}
\Pturn{Compute
\(\displaystyle
 \pi_{\rm open}\leftarrow
 \mathsf{RecProve}(\mathsf{pp},\rho,B,\beta,P;\tr).
\)}
\Psend{\Splitbadge{\pi_{\rm open}}}
\Vturn{Require
\(\begin{aligned}
 &\mathsf{RecVerify}
   (\mathsf{pp},\rho,B,\beta;\tr,c_N,c_H)=1,\\
 &\mathsf{NoPendingHint}=1,\\
 &\mathsf{NargEOF}(c_N)=\mathsf{HintEOF}(c_H)=1,\\
 &\mathsf{HostEOF}(\pi_{\rm host})=1,\\
 &\mathsf{Counts}(c_N,c_H)=(n_N,n_H).
\end{aligned}\)}
\end{pvlongexchange}
\end{pvstep}

\setlength{\textwidth}{7.45in}
\setlength{\columnwidth}{7.45in}
\setlength{\linewidth}{7.45in}
\setlength{\hsize}{7.45in}

P9/V9 follows the recursive-opening ledger in
Section~\ref{sec:wire-recursive-ledger}.  Its proof spans both streams; the
whole $\pi_{\rm open}$ \mustnot{} be placed on the hint channel.

\section{Verifier algorithm}
\label{sec:verifier}

\begingroup\footnotesize

\definecolor{verifierNarg}{HTML}{315B7D}
\definecolor{verifierHint}{HTML}{A15C00}
\definecolor{verifierAbsorb}{HTML}{2D6F67}
\definecolor{verifierSqueeze}{HTML}{6D4CC2}
\definecolor{verifierCheck}{HTML}{A81E1E}
\definecolor{verifierDerived}{HTML}{667085}
\definecolor{verifierCall}{HTML}{1B7F3B}

\newcommand{\VBadge}[2]{%
  \textcolor{#1}{\text{\scriptsize\sffamily\bfseries[#2]}}}
\newcommand{\VNargOp}[1]{%
  \VBadge{verifierNarg}{N}\,\textcolor{verifierNarg}{#1}}
\newcommand{\VHintOp}[1]{%
  \VBadge{verifierHint}{H}\,\textcolor{verifierHint}{#1}}
\newcommand{\VAbsorbOp}{%
  \VBadge{verifierAbsorb}{A}\,\textcolor{verifierAbsorb}{\Absorb}}
\newcommand{\VSqueezeOp}{%
  \VBadge{verifierSqueeze}{S}\,\textcolor{verifierSqueeze}{\Squeeze}}
\newcommand{\VCheckOp}{%
  \VBadge{verifierCheck}{C}\,\textcolor{verifierCheck}{\mathsf{Require}}}
\newcommand{\VDerive}{\VBadge{verifierDerived}{D}\,}
\newcommand{\VCall}[1]{%
  \VBadge{verifierCall}{CALL}\,\textcolor{verifierCall}{\mathsf{#1}}}
\newcommand{\VReadNarg}{\VNargOp{\mathsf{ReadNarg}}}
\newcommand{\VReadHint}{\VHintOp{\mathsf{ReadHint}}}
\newcommand{\PVStep}[2]{%
  \textbf{\textcolor{specnavy}{P#1/V#1: #2}}}

The verifier rejects on a failed gate, malformed value, unsupported public
configuration, or cursor mismatch.  NARG and hint cursors are independent.  A
hint remains pending until its predicate succeeds, and every squeeze requires
$\mathsf{NoPendingHint}$.

\begingroup\scriptsize
\setlength{\tabcolsep}{2pt}
\renewcommand{\arraystretch}{1.02}
\begin{tabularx}{\linewidth}{@{}*{7}{>{\centering\arraybackslash}X}@{}}
\rowcolor{specpale}
\textcolor{verifierNarg}{\textbf{[N]}} read / bind&
\textcolor{verifierHint}{\textbf{[H]}} read / check&
\textcolor{verifierAbsorb}{\textbf{[A]}} absorb&
\textcolor{verifierSqueeze}{\textbf{[S]}} challenge&
\textcolor{verifierCheck}{\textbf{[C]}} reject gate&
\textcolor{verifierDerived}{\textbf{[D]}} local value&
\textcolor{verifierCall}{\textbf{[CALL]}} subprotocol\\[-1pt]
\end{tabularx}
\begin{tabularx}{\linewidth}{@{}*{4}{>{\centering\arraybackslash}X}@{}}
\textbf{P1/V1--P2/V2} statement / PIOP&
\textbf{P3/V3} input Sumcheck&
\textbf{P4/V4--P6/V6} virtual-witness core&
\textbf{P7/V7--P9/V9} adjoint / opening\\
\end{tabularx}
\endgroup

P1/V1 and P2/V2 belong to the configured end-to-end wrapper.  The F2Z
verifier below begins at P3/V3 after the PIOP verifier has produced
$(\bm v,\mu)$.

For a canonical domain byte string $D$ and canonical payload $z$, the notation
$\VAbsorbOp(D;z)$ abbreviates the frame-free operation
\[
 \Absorb\!\left(D\parallel\LE{64}{\ByteLen(z)}\parallel z\right)
\]
from Equation~\eqref{eq:transcript-absorb}.  A read returns the decoded value
and its exact canonical payload bytes; it never returns a frame object.  The
named domains below are the verifier-derived encodings registered in
Section~\ref{sec:wire-events}.

\paragraph{Procedure $\mathsf{VerifyEndToEnd}$.}
\begin{algorithmic}[1]
\Require $\mathsf{pp},x,\mathsf{com}=(\rho),
  \pi_{\rm e2e}=(\pi_{\rm PIOP},\pi_{\rm host})$
\Ensure $\mathsf{true}$ or $\mathsf{false}$

\Statex
\Statex \PVStep{1}{Bind the constraint statement and commitment}
\State $\VCheckOp[\mathsf{ValidateStatementConfigRoot}
  (\mathsf{pp},x,\rho)=1]$

\Statex
\Statex \PVStep{2}{Run the configured PIOP verifier}
\State $(\mathsf{ok}_{\rm PIOP},\bm v,\mu)\gets
  \VCall{PIOP.Verify}(\mathsf{pp}_{\rm PIOP},x,\rho,\pi_{\rm PIOP})$
\State $\VCheckOp[\mathsf{ok}_{\rm PIOP}=1
  \land |\bm v|=n_h\land\mu\in\Fq]$
\State \Return $\mathsf{VerifyF2Z}
  (\mathsf{pp},(\bm v,\mu,\mathsf{configuration}),
   (\rho),\pi_{\rm host})$
\end{algorithmic}

\paragraph{Procedure $\mathsf{VerifyF2Z}$.}
\begin{algorithmic}[1]
\Require $\mathsf{pp},\mathcal I_{\rm LinBitsDual}=(\bm v,\mu,\mathsf{configuration}),
  \mathsf{com}=(\rho),\pi_{\rm host}$
\Ensure $\mathsf{true}$ or $\mathsf{false}$

\Statex
\Statex \PVStep{3}{Reduce the PIOP handoff with Sumcheck}
\Statex \hspace{\algorithmicindent}
  $\mathcal R_{\rm LinBitsDual}[\bm h]$
\State $(N_{\rm F2Z},H_{\rm F2Z},n_N,n_H,b_N,b_H)
  \gets\mathsf{DecodeHostV1}(\pi_{\rm host})$
\State $\VCheckOp[\mathsf{HeaderV1}=1\land\mathsf{wire\_version}=1]$
\State $\VCheckOp[\mathsf{ValidateConfig}
  (\mathsf{pp},\mathcal I_{\rm LinBitsDual},\rho)=1]$
\State Read $(n_f,m_f,n_h,m_h,d_{F,r},d_{F,c},d_r,d_c,g)$ from the fixed public
  configuration
\State $\VCheckOp[n_f=2^{m_f}\land n_h=2^{m_h}\land m_h=d_r+d_c
  \land 128\mid n_f\land m_f=d_{F,r}+d_{F,c}
  \land m_f\ge7\land 2^{d_r}q<2^{127}]$
\State $\VDerive m_p\gets m_f-7$
\State $\VDerive b_{\rm virt}\gets
  \mathsf{HasVirtualization}(\mathsf{configuration})$
\If{$b_{\rm virt}$}
  \State Read $\bm M$ from the developer configuration and
    $\VCheckOp[\bm M\in\Ftwo^{n_h\times n_f}]$
\Else
  \State $\VCheckOp[n_h=n_f]$ and
         $\VCheckOp[\mathsf{SameFlatOrder}(\mathsf{configuration})=1]$
  \State $\VDerive\bm M\gets I_{n_f}$
\EndIf
\State $\VCheckOp[\operatorname{ord}(g)=2^{128}-1]$
\State $\VCheckOp[|\bm v|=n_h
  \land\mu\in\Fq]$
\State $\VDerive\tr\gets\tr_{\mathsf{init}}$
\State $\VAbsorbOp(\mathsf D_{\rm domain};\mathsf{DomainV1})$
  using Equation~\eqref{eq:domain-statement-payloads}
\State $\VAbsorbOp(\mathsf D_{\rm statement};\mathsf{StatementV1})$
  using Equation~\eqref{eq:domain-statement-payloads}
\State Initialize independent cursors
  $\VDerive c_N\gets\mathsf{Cursor}(N_{\rm F2Z},n_N,b_N)$,
  $c_H\gets\mathsf{Cursor}(H_{\rm F2Z},n_H,b_H)$, and
  $\mathsf{PendingHint}\gets0$
\State Load the fixed input-Sumcheck, GKR, row-LinCheck, and opening ledgers

\State $\VDerive\sigma_{-1}^{\rm shape}\gets\mu$
\For{$i=0,\ldots,m_h-1$}
  \State $(g_i^{\rm shape},p_i^{\rm shape})\gets
    \VReadNarg(c_N,\mathsf D_{{\rm shape},i})$
  \State $\VCheckOp[g_i^{\rm shape}\in\Fq[X]\land
    \deg g_i^{\rm shape}\le2\land
    g_i^{\rm shape}(0)+g_i^{\rm shape}(1)=
      \sigma_{i-1}^{\rm shape}]$
  \State $\VAbsorbOp(\mathsf D_{{\rm shape},i};p_i^{\rm shape})$
  \State $\VDerive(\tr,s_{h,i})\gets
    \mathsf{SqueezeFqV1}(\tr,\mathsf D_{{\rm shape\text{-}challenge},i})$
  \State $\VDerive\sigma_i^{\rm shape}\gets
    g_i^{\rm shape}(s_{h,i})$
\EndFor
\State $\VDerive\bm s_h\gets(s_{h,0},\ldots,s_{h,m_h-1})$,
  $\tau\gets\sigma_{m_h-1}^{\rm shape}$,
  $\kappa\gets\widetilde{\bm v}(\bm s_h)$
\State Split $\bm s_h=(\bm s_r,\bm s_c)$ with
  $|\bm s_r|=d_r$ and $|\bm s_c|=d_c$
\State $\VDerive a_r\gets\eq_{\Fq}(r,\bm s_r)$ and
  $a'_c\gets\kappa\eq_{\Fq}(c,\bm s_c)$
\State $\VCheckOp[|\bm a|=2^{d_r}\land|\bm a'|=2^{d_c}
  \land\sum_ra_r=1]$
\State Retain the terminal claim
  $\kappa\widetilde H_q(\bm s_h)=\tau$
\State $\VAbsorbOp(\mathsf D_{\rm seam};\mathsf{LinCheckSeamV1})$
  using Equation~\eqref{eq:lincheck-seam-payload}
\State For every row $r$, derive the unique
  $\gamma_r\in[0,q)$ with $\pi_q(\gamma_r)=a_r$
\State Fix, without advancing either cursor, the remaining record schedule

\Statex
\Statex \PVStep{4}{Fold each column and encode the folds as product claims}
\State $(\bm\nu,p_\nu)\gets
  \VReadNarg(c_N,\mathsf D_{\nu})$
  \State $\VCheckOp[|\bm\nu|=2^{d_c}\land\bm\nu\in(\K^\times)^{2^{d_c}}]$
  \State $\VAbsorbOp(\mathsf D_{\nu};p_\nu)$
  \State $\mathsf{PendingHint}\gets1$,
    $(\bm\mu,p_\mu)\gets
      \VReadHint(c_H,\mathsf D_{\mu})$
  \State $\VCheckOp[|\bm\mu|=2^{d_c}\land
    \forall c:\ 0\le\mu_c<2^{127}\land g^{\mu_c}=\nu_c]$
  \State $\VCheckOp[\sum_ca'_c\pi_q(\mu_c)=\tau]$
  \State $\mathsf{PendingHint}\gets0$
  \State $(\tr,\zeta)\gets
    \VSqueezeOp_{\K}^{d_c}(\tr;\mathsf D_\zeta)$,
    $e_0\gets\widetilde{\bm\nu}(\zeta)$

  \Statex
  \Statex \PVStep{5}{Reduce PESAT to a weighted sum with GKR}
  \Statex \hspace{\algorithmicindent}
    $\mathcal R_{\rm PESAT}[\bm h]
      \xRightarrow{\Pi_{\rm GKR}}\mathcal R_{\rm sum}^{h}$
  \State $\VDerive\mathcal I_{\rm GKR}\gets
    \mathsf{GKRInput}(\bm\gamma,\bm\nu,\zeta,e_0)$
  \State $(\mathsf{ok}_{\rm GKR},
    \tr',\bm\omega_\xi,\sigma)
    \gets\VCall{GKR.Verify}
    (\mathsf{pp}_{\rm GKR},\mathcal I_{\rm GKR};
     \tr,c_N,c_H)$
  \State $\VCheckOp[\mathsf{ok}_{\rm GKR}=1]$
  \State $\VDerive\tr\gets\tr'$

  \Statex
  \Statex \PVStep{6}{Reduce the weighted sum to an MLE claim}
  \Statex \hspace{\algorithmicindent}
    $\mathcal R_{\rm sum}^{h}
      \xRightarrow{\Pi_{\rm MLE}}\mathcal R_{\rm eval}^{h}$
  \State $\VDerive\mathcal I_{\rm MLE}\gets
    (\bm\omega_\xi,\sigma,\mathsf{configuration})$
  \State $(\mathsf{ok}_{\rm MLE},
    \tr',\bm r_h,\theta,\eta)
    \gets\VCall{RowLinCheck.Verify}
    (\mathcal I_{\rm MLE};\tr,c_N,c_H)$
  \State $\VCheckOp[\mathsf{ok}_{\rm MLE}=1]$ and
    $\VDerive\tr\gets\tr'$
  \State $\VDerive\lambda\gets
    \widetilde{\bm\omega_\xi}(\bm r_h)$
  \State $\VCheckOp[\theta=\lambda\eta]$
  \Statex \hspace{\algorithmicindent}
    No absorb or squeeze occurs at the semantic P5/V5--P6/V6 boundary;
    the fixed record contains $(\pi_{\rm GKR},\pi_{\rm MLE})$ in that order.

\Statex \PVStep{7}{Rewrite $\langle\bm a,\bm M\bm f\rangle$ through the adjoint}
\State $\VDerive\bm a[z]\gets\eq(z,\bm r_h)$
\Statex \hspace{\algorithmicindent}
  $\langle\bm a,H_{\K}\rangle_{\K}
   =\langle\bm a,\bm M_{\K}\bm f_{\K}\rangle_{\K}
   =\langle\bm M_{\K}^{\mathsf T}\bm a,\bm f_{\K}\rangle_{\K}$
\State $\VDerive\bm a_f\gets\bm M_{\K}^{\mathsf T}\bm a$
\State $\VDerive A_f[v,j]\gets\bm a_f[128j+v]$

\Statex \PVStep{8}{Ring-switch the adjoint claim into the committed packing}
\State $\VDerive
  (\Theta,(A_\theta,B_\theta)_{\theta\in\Theta})
  \gets\VCall{RingDecomp}(\bm a_f)$
\State $\VCheckOp[\Theta\ne\varnothing\land
  \forall v,j:\ A_f[v,j]
    =\sum_{\theta\in\Theta}A_\theta[v]B_\theta[j]]$
\For{$\theta\in\Theta$ in canonical order}
  \State $(\bm s_\theta,p_{s,\theta})\gets
    \VReadNarg(c_N,\mathsf D_{s,\theta})$
  \State $\VCheckOp[|\bm s_\theta|=128]$
  \State $\VAbsorbOp(\mathsf D_{s,\theta};p_{s,\theta})$
\EndFor
\State $\VCheckOp[
  \sum_{\theta\in\Theta}\sum_{v=0}^{127}
    A_\theta[v]s_{\theta,v}=\eta]$
\State $(\tr,r'')\gets\VSqueezeOp_{\K}^7(\tr;\mathsf D_{r''})$
\State $\VDerive\Phi_{r''}(x)\gets
  \sum_u\eq(u,r'')[x]_u$
\For{$\theta\in\Theta$ in canonical order}
  \State $\VDerive t_{\theta,u}\gets
    \sum_v[s_{\theta,v}]_u\beta_v$
  \State $\VDerive\beta_\theta\gets
    \sum_u\eq(u,r'')t_{\theta,u}$,
    $\bm\omega_\theta[j]\gets\Phi_{r''}(B_\theta[j])$
  \State $(\tr,\lambda_\theta)\gets
    \VSqueezeOp_{\K}(\tr;\mathsf D_{\lambda,\theta})$
\EndFor
\State $\VDerive B\gets
  \sum_{\theta\in\Theta}\lambda_\theta\bm\omega_\theta$,
  $\beta\gets\sum_{\theta\in\Theta}\lambda_\theta\beta_\theta$

\Statex \PVStep{9}{Open the packed claim}
\State $\VDerive p_{\rm target}\gets
  \LE{32}{m_p}\parallel\operatorname{Enc}_{\K}(\beta)$
\State $\VAbsorbOp(\mathsf D_{\rm target};p_{\rm target})$
\State $\VCheckOp[\VCall{RecVerify}
  (\mathsf{pp},\rho,B,\beta;\tr,c_N,c_H)=1]$
\State $\VCheckOp[\mathsf{NoPendingHint}=1\land
  \mathsf{NargEOF}(c_N)=1\land\mathsf{HintEOF}(c_H)=1\land
  \mathsf{HostEOF}(\pi_{\rm host})=1\land
  \mathsf{Counts}(c_N,c_H)=(n_N,n_H)]$
\State \Return $\mathsf{true}$
\end{algorithmic}

The helper verifiers own their subordinate ledgers.  Every NARG read advances
only $c_N$, every hint read advances only $c_H$, and no helper returns with an
unchecked hint.  P4/V4's $\Nbadge{\nu}$, $\Hbadge{\mu}$ order still uses
the NARG-then-Hint order fixed in Section~\ref{sec:interaction}.
The byte rules in Section~\ref{sec:wire-format} govern every read, absorb,
squeeze, PoW check, and EOF condition.

\endgroup

\section{Wire and transcript format v1 (Optional)}
\label{sec:wire-format}

\begingroup\small

\begin{normative}
This optional section defines the concrete v1 profile.  Optionality concerns
whether an implementation selects this profile; once selected, every rule in
this section is mandatory.  The section fixes the bytes consumed by a
conforming v1 prover and verifier.
The proof system has two physical, canonically encoded positional record streams, a
BLAKE3 transcript, SHA-256 Merkle roots, and the recursive Ligerito and
proof-of-work parameters below.  A proof using the legacy positional/bincode
codec is not a v1 proof.
\end{normative}

\subsection{Version and fixed algorithms}
\label{sec:wire-version}

The eight-byte container magic is the ASCII string \code{F2ZPCS\textbackslash0\textbackslash0},
and the wire version is one.  This document defines the not-yet-frozen v1
format; its input-shaping ledger supersedes the earlier draft that reserved
the entire \code{1000--1fff} range.  Once a complete v1 proof/transcript
fixture is published, changing any normative rule in this section requires a
new wire version.

The transcript and its XOF are unkeyed BLAKE3; Merkle nodes and roots use
SHA-256; PoW uses unkeyed BLAKE3-256.  The field basis is the polynomial basis
fixed in Section~\ref{sec:domains}.  The public parameters fix the
GKR protocol, constrained-code construction, transcript, and opening protocol.
None is selected or negotiated by proof bytes.  A proof system without a
byte-exact GKR ledger or without a canonical codeword construction and root
vectors is incomplete and \mustnot{} accept proofs.

Write $m:=m_f$ and $m_p:=m-7$ for the packed-vector dimension.  For the deployed
range $22\le m\le35$, let $J_m$ be the number of recursive
opening levels.  The level shapes and all parameters except the base fold
difficulties are fixed by
\begin{equation}
 \begin{aligned}
  h_i&:=m-11-3i,&
  J_m&:=\min\{J\ge1:h_{J-1}\in\{3,4,5\}\},\\
  r_i&:=1+i,&
  k_i&:=\begin{cases}4,&i=0,\\3,&i>0,\end{cases}\\
  Q_i&:=(183,90,60,45,36,31,27,24)_i,&
  g_i&:=16,\\
  o_i&:=\begin{cases}0,&i=0,\\1,&i>0.\end{cases}
 \end{aligned}
 \label{eq:ligerito-v1-shape}
\end{equation}
Here $h_i$ is the message-column exponent, $r_i$ the inverse-rate exponent,
$k_i$ the fold count, $Q_i$ the query count, $g_i$ the query-grinding
difficulty, and $o_i$ the OOD sample count.  Implementations \must{} use the
integers in Equation~\eqref{eq:ligerito-v1-shape}, not a floating-point
recomputation.

The remaining base fold difficulties $f_i$ are the following literal table.
Only the first $J_m$ entries of the row for $m$ exist.

\begingroup\footnotesize
\renewcommand{\arraystretch}{1.02}
\begin{longtable}{@{}c@{\quad}l@{}}
\toprule
$m$&$(f_0,\ldots,f_{J_m-1})$\\
\midrule
22&$(9,6,3)$\\
23&$(10,7,4,1)$\\
24&$(11,8,5,2)$\\
25&$(12,9,6,3)$\\
26&$(13,10,7,4,2)$\\
27&$(14,11,8,5,3)$\\
28&$(15,12,9,6,4)$\\
29&$(16,13,10,7,5,3)$\\
30&$(17,14,11,8,6,4)$\\
31&$(18,15,12,9,7,5)$\\
32&$(19,16,13,10,8,6,3)$\\
33&$(20,17,14,11,9,7,4)$\\
34&$(21,18,15,12,10,8,5)$\\
35&$(22,19,16,13,11,9,6,6)$\\
\bottomrule
\end{longtable}
\endgroup

\subsection{Primitive canonical encodings}
\label{sec:wire-primitives}

All unsigned integers are fixed-width little-endian.  Decoders \mustnot{}
reduce, truncate, ignore high bits, or convert to a host-sized integer before
checking that the value fits.  The following encodings are the only canonical
ones.

\begingroup\footnotesize
\renewcommand{\arraystretch}{1.05}
\begin{tabularx}{\linewidth}{@{}p{1.25in}p{1.35in}X@{}}
\toprule
\textbf{Type}&\textbf{Bytes}&\textbf{Canonicality rule}\\
\midrule
$u_8,u_{16},u_{32},u_{64},u_{128}$&1, 2, 4, 8, 16&Unsigned
little-endian of exactly the named width.\\
$a\in\K$&16&$\LE{64}{a_{\lo}}\parallel\LE{64}{a_{\hi}}$ in the fixed
polynomial basis; coefficient bit $j$ is bit $j$ of the resulting 128-bit
word.\\
$a\in\Fq$&16&$\LE{128}{\bar a}$, where $0\le\bar a<q$ and
$\pi_q(\bar a)=a$; reject an integer at least $q=2^{100}-15$.\\
Bounded exponent $u$&16&$\LE{128}{u}$; require $u<2^{127}$.\\
Nonce&8&$\LE{64}{n}$.\\
Merkle value&32&Raw SHA-256 digest bytes, in digest output order.\\
Vector $[x_0,\ldots,x_{n-1}]$&4 plus elements&$\LE{32}{n}$ followed by
canonical elements in index order; $n$ must equal the protocol-derived expected
count.\\
Opened rows and path&nested&Vector of rows, each a vector of $\K$ elements,
followed by a vector of 32-byte Merkle values.  Expected row, width, and path
counts are derived before allocation.\\
\bottomrule
\end{tabularx}
\endgroup

Every composite is concatenation without padding.  Options, Rust enums,
\code{usize}, varints, serde defaults, and bincode layouts are not wire types.
For every accepted primitive or composite $x$, decoding its encoding must
consume the whole bounded slice and re-encoding the result must reproduce the
same bytes.

\subsection{Merkle rows and canonical multiproofs}
\label{sec:wire-merkle}

For opening level $0\le k<J_m$, define the message dimension, tree depth,
leaf count, and row width by
\begin{equation}
 n_k:=h_k+k_k,\qquad D_k:=h_k+r_k,\qquad
 \Lambda_k:=2^{D_k},\qquad I_k:=2^{k_k}.
 \label{eq:merkle-level-geometry}
\end{equation}
Equation~\eqref{eq:merkle-level-geometry} implies
$n_0=m-7$, $n_k=h_{k-1}$ for $k>0$, $D_k=m-10-2k$,
$I_0=16$, and $I_k=8$ for $k>0$.  Let
$C_k\in\K^{\Lambda_k\times I_k}$ be the level-$k$ codeword matrix, with
leaves and lanes in increasing integer order.  Its exact row and tree bytes
are
\begin{equation}
 \begin{aligned}
  \mathsf{RowBytes}_k(q)
   &:=\operatorname{Enc}_{\K}(C_k[q,0])\parallel\cdots\parallel
      \operatorname{Enc}_{\K}(C_k[q,I_k-1]),\\
  H_{k,0,q}&:=\operatorname{SHA256}(\mathsf{RowBytes}_k(q)),\\
  H_{k,d+1,j}&:=\operatorname{SHA256}
      (H_{k,d,2j}\parallel H_{k,d,2j+1}),\\
  \mathsf{root}_k&:=H_{k,D_k,0}.
 \end{aligned}
 \label{eq:merkle-v1-tree}
\end{equation}
Here $\operatorname{Enc}_{\K}$ is the 16-byte encoding in
Section~\ref{sec:wire-primitives}.  Equation~\eqref{eq:merkle-v1-tree} uses
no native-memory cast, row-length prefix, row index, level domain code, padding, or
leaf/internal-node domain byte.  Adding such a byte changes every root and
requires a new wire version.  The fixed code specification defines
how $C_k$ is generated; this section fixes its portable row and tree encoding.

Let the already-derived query set be
$A_0=(q_0<\cdots<q_{Q_k-1})$.  At depth $d$, scan the active positions $A_d$
in increasing order.  If both siblings are active, combine the pair once and
emit no digest.  Otherwise emit the digest at position
$p\mathbin{\mathsf{xor}}1$, using $p$'s parity to place it on the left or
right.  Then set
$A_{d+1}:=\operatorname{sortuniq}\{\lfloor p/2\rfloor:p\in A_d\}$.
The canonical multiproof is the concatenation of emitted digests, bottom-up
and in that scan order, and its exact count is
\begin{equation}
 M_k:=\sum_{d=0}^{D_k-1}
 \left|\{p\in A_d:(p\mathbin{\mathsf{xor}}1)\notin A_d\}\right|.
 \label{eq:merkle-proof-count}
\end{equation}
Equation~\eqref{eq:merkle-proof-count} fixes the exact number of transmitted
path digests.  The level-$k$ opening payload is exactly
\begin{equation}
 \LE{32}{Q_k}\parallel
 \prod_{j=0}^{Q_k-1}
   \bigl(\LE{32}{I_k}\parallel\mathsf{RowBytes}_k(q_j)\bigr)
 \parallel\LE{32}{M_k}\parallel
 \prod_{u=0}^{M_k-1}\mathsf{path}_k[u],
 \label{eq:merkle-opening-payload}
\end{equation}
with byte length
\begin{equation}
 8+Q_k(4+16I_k)+32M_k.
 \label{eq:merkle-opening-length}
\end{equation}
The verifier derives $A_d$, $M_k$, and
Equation~\eqref{eq:merkle-opening-length} from the sampled queries before
allocating or reading the bounded Hint record, requires rows in $q_j$ order,
reconstructs Equation~\eqref{eq:merkle-v1-tree}, and rejects unused path
digests or payload bytes.

\subsection{Outer two-stream container}
\label{sec:wire-container}

Here $N$ and $H$ are the concatenations of their canonical payload records in
the verifier-event and hint-check orders, respectively.  They contain no
transmitted event domain, scope, or generic record-length wrapper; each next
decoder and its bounded length follow from the fixed ledger and the canonical
payload type.

The 40-byte header is fixed as follows.  Offsets and sizes are decimal; all
multibyte integers use Section~\ref{sec:wire-primitives}.

\begin{lstlisting}[style=wire]
offset  size  field
0       8     magic = 46 32 5a 50 43 53 00 00
8       2     wire_version = 01 00
10      2     header_len = 28 00                 (40)
12      4     flags = 00 00 00 00
16      4     narg_record_count
20      4     hint_record_count
24      8     narg_byte_len
32      8     hint_byte_len
40      ...   NARG records, then Hint records
\end{lstlisting}

Writing $\mathsf{Hdr}_{40}$ for those exact header bytes, the only v1 host
encoding is
\begin{equation}
 \HostEncode(N,H)
 :=\mathsf{Hdr}_{40}\parallel N\parallel H,
 \qquad
 \ByteLen(\HostEncode(N,H))=40+\ByteLen(N)+\ByteLen(H).
 \label{eq:host-encoding-v1}
\end{equation}
Equation~\eqref{eq:host-encoding-v1} implies the four benchmark metrics
\begin{equation}
 \begin{aligned}
  \NargLen(\pi)&=\ByteLen(N),&
  \HintLen(\pi)&=\ByteLen(H),\\
  \PayloadLen(\pi)&=\ByteLen(N)+\ByteLen(H),&
  \WireLen(\pi)&=40+\PayloadLen(\pi).
 \end{aligned}
 \label{eq:wire-lengths-v1}
\end{equation}
The decoder \must{} verify Equation~\eqref{eq:wire-lengths-v1} with checked
arithmetic before allocating, enforce statement-derived upper bounds, and
require input length exactly $40+\mathtt{narg\_byte\_len}+
\mathtt{hint\_byte\_len}$.  The two declared record counts must match the
records actually consumed.  Header count and length fields are transport
metadata and are not transcript input.

\subsection{Typed event records and domain registry}
\label{sec:wire-events}

Every absorbed NARG value, public value, and squeeze event uses the three-part
encoding from Equation~\eqref{eq:transcript-absorb}.  Its canonical domain is
the following 16-byte string:

\begin{equation}
 D_{\tau,j,k,i}:=
 \LE{16}{\tau}\parallel\LE{16}{0}\parallel
 \LE{32}{j}\parallel\LE{32}{k}\parallel\LE{32}{i}.
 \label{eq:event-domain-v1}
\end{equation}

Here $\tau$ is the numeric semantic-domain code and $(j,k,i)$ is its
verifier-derived scope.  Let $\bot_{32}:=\mathtt{0xffffffff}$ denote an
inapplicable coordinate.  These values are fixed by the current ledger
position; they are never read from either proof stream.

The names used by the verifier algorithm in Section~\ref{sec:verifier}
abbreviate
\begin{equation}
 \begin{aligned}
  \mathsf D_{\rm domain}&:=D_{\mathtt{0x0001},\bot_{32},\bot_{32},\bot_{32}},&
  \mathsf D_{\rm statement}&:=D_{\mathtt{0x0002},\bot_{32},\bot_{32},\bot_{32}},\\
  \mathsf D_{{\rm shape},i}&:=D_{\mathtt{0x1002},\bot_{32},\bot_{32},i},&
  \mathsf D_{{\rm shape\text{-}challenge},i}&:=
    D_{\mathtt{0x1101},\bot_{32},\bot_{32},i},\\
  \mathsf D_{\rm seam}&:=D_{\mathtt{0x1003},\bot_{32},\bot_{32},\bot_{32}},&
  \mathsf D_{\nu}&:=D_{\mathtt{0x2001},\bot_{32},\bot_{32},\bot_{32}},\\
  \mathsf D_{\mu}&:=D_{\mathtt{0x2081},\bot_{32},\bot_{32},\bot_{32}},&
  \mathsf D_{s,\theta}&:=D_{\mathtt{0x4001},\theta,\bot_{32},\bot_{32}},&
  \mathsf D_{\lambda,\theta}&:=D_{\mathtt{0x4102},\theta,\bot_{32},\bot_{32}},\\
  \mathsf D_{\rm target}&:=D_{\mathtt{0x5001},\bot_{32},\bot_{32},\bot_{32}}.&&
 \end{aligned}
 \label{eq:verifier-domain-abbreviations}
\end{equation}
The vector names $\mathsf D_\zeta$ and $\mathsf D_{r''}$ denote consecutive
coordinate domains $D_{\mathtt{0x2101},\bot_{32},\bot_{32},i}$ and
$D_{\mathtt{0x4101},\bot_{32},\bot_{32},i}$, respectively.

For payload $z$, an absorbed event appends exactly
\begin{equation}
 D_{\tau,j,k,i}
 \parallel\LE{64}{\ByteLen(z)}\parallel z.
 \label{eq:event-absorb-v1}
\end{equation}
There is no frame or event wrapper.  A transmitted NARG or Hint record is only
the canonical payload $z$.  Its decoder and expected length are determined by
the fixed event ledger, including any canonical length prefix belonging to the
payload type itself.  Thus proof parsing is positional and separate from
transcript absorption: a NARG read decodes and validates $z$, then applies
Equation~\eqref{eq:event-absorb-v1}; a Hint read decodes and checks $z$ but
never applies that equation.

The domain registry is below.  For \textsf{D}, \textsf{N}, and \textsf{S}
events, domain codes and scopes are verifier-derived transcript inputs; for
\textsf{H} records they are verifier-derived ledger selectors.  None is a
record header.  \textsf{D} events are public or derived, have no proof record,
and are absorbed at their ledger position; \textsf{N} payloads occur only in
$N$ and are absorbed; \textsf{H} payloads occur only in $H$ and are checked but never absorbed;
\textsf{S} domains are squeeze requests.  For an \textsf{S} base domain code $\tau$, the
result domain code is $\tau\mathbin{\mathrm{OR}}\mathtt{0x8000}$ and is derived,
never transmitted.  For an \textsf{S} row, the last column describes the
derived result; its request payload is always $\LE{32}{n}\parallel a$ as fixed
by Equation~\eqref{eq:transcript-v1}.

\begingroup\scriptsize
\renewcommand{\arraystretch}{1.04}
\setlength{\tabcolsep}{2pt}
\begin{longtable}{@{}p{0.52in}p{0.30in}p{1.62in}p{4.72in}@{}}
\toprule
\textbf{Code}&\textbf{Class}&\textbf{Event}&\textbf{Scope; proof payload or squeeze result}\\
\midrule
\endfirsthead
\toprule
\textbf{Code}&\textbf{Class}&\textbf{Event}&\textbf{Scope; proof payload or squeeze result}\\
\midrule
\endhead
\code{0001}&D&domain/version&$(\bot,\bot,\bot)$; $\mathsf{DomainV1}$ from
Equation~\eqref{eq:domain-statement-payloads}.\\
\code{0002}&D&statement&$(\bot,\bot,\bot)$; $\mathsf{StatementV1}$ from
Equation~\eqref{eq:domain-statement-payloads}.\\
\code{0003}&---&legacy core seam&Reserved and \mustnot{} be scheduled; the input
Sumcheck terminates with derived domain code \code{1003}.\\
\code{1001}&---&PIOP row collapse&Reserved for a legacy PIOP-side phase;
it does not occur when F2Z receives the public $\bm v$-weighted claim.\\
\code{1002}&N&input-Sumcheck round&$(\bot,\bot,i)$ for $0\le i<m_h$; exactly
$\operatorname{Enc}_{\Fq}(g_i(0))\parallel
 \operatorname{Enc}_{\Fq}(g_i(1))\parallel
 \operatorname{Enc}_{\Fq}(g_i(2))$ for the degree-two round polynomial, with
no vector prefix.\\
\code{1003}&D&input-Sumcheck seam&$(\bot,\bot,\bot)$;
$\mathsf{LinCheckSeamV1}$ from
Equation~\eqref{eq:lincheck-seam-payload}.\\
\code{1101}&S&input-Sumcheck challenge&$(\bot,\bot,i)$ for $0\le i<m_h$;
one $\Fq$ challenge from the rejection sampler below.\\
\code{1000}&---&reserved&The verifier \mustnot{} schedule this domain code.\\
\code{1004--1100}&---&reserved&The verifier \mustnot{} schedule these domain codes.\\
\code{1102--1fff}&---&reserved&The verifier \mustnot{} schedule these domain codes.\\
\code{2001}&N&fold exponent image $\nu$&$(\bot,\bot,\bot)$; vector of
$2^{d_c}$ elements of $\K$.\\
\code{2081}&H&integer fold $\mu$&$(\bot,\bot,\bot)$; vector of
$2^{d_c}$ canonical $u_{128}<2^{127}$.\\
\code{2101}&S&fold challenge $\zeta$&$(\bot,\bot,i)$ for $0\le i<d_c$; one $\K$
challenge.\\
\code{3001}&N&GKR phase-$c$ sum&$(\bot,k,\bot)$; one $\K$ element, absorbed
after the phase-link check and before the first phase-$c$ round.\\
\code{3002--30ff}&Sub.&subprotocol messages&Except for reserved code
\code{3081}, the fixed subordinate ledger owns exact message domain codes,
codecs, and checks for $\pi_{\rm GKR}$ followed by $\pi_{\rm MLE}$.\\
\code{3100--31ff}&Sub.&subprotocol squeezes&The fixed subordinate ledger
owns exact squeeze domain codes and result types in the same order.\\
\code{4001}&N&ring $\bm s_\theta$ vector&$(\theta,\bot,\bot)$ for
$\theta\in\Theta$ in canonical order; vector of exactly 128 elements of $\K$.\\
\code{4002}&---&reserved&The verifier \mustnot{} schedule this domain code.\\
\code{4101}&S&shared $r''$&$(\bot,\bot,i)$ for $0\le i<7$; one $\K$
challenge.\\
\code{4102}&S&ring batch coefficient $\lambda_\theta$&$(\theta,\bot,\bot)$
for $\theta\in\Theta$ in canonical order; one $\K$ challenge.\\
\code{4103}&---&reserved&The verifier \mustnot{} schedule this domain code.\\
\code{5001}&D&opening target&$(\bot,\bot,\bot)$;
$\LE{32}{m_p}\parallel\operatorname{Enc}_{\K}(\beta)$.  $B$ is recomputed.\\
\code{5002}&N&initial root&$(\bot,0,\bot)$; one 32-byte SHA-256 digest.\\
\code{5003}&N&recursive root&$(\bot,k,\bot)$, $1\le k<J_m$; one digest.\\
\code{5004}&N&sumcheck message&scope by kind below; $\mathsf{kind}:u_8$
followed by $\operatorname{Enc}_{\K}(u_0)\parallel
\operatorname{Enc}_{\K}(u_2)$.\\
\code{5005}&N&OOD value&$(a,k,\bot)$; one $\K$ element for sample $a$.\\
\code{5006}&N&final $y_r$&$(\bot,J_m-1,\bot)$; vector of exactly
$2^{h_{J_m-1}}$ elements of $\K$.\\
\code{5007}&N&fold-PoW nonce&$(\bot,k,i)$; one $u_{64}$.\\
\code{5008}&N&query-PoW nonce&$(\bot,k,\bot)$; one $u_{64}$.\\
\code{5080}&H&initial opening&$(\bot,0,\bot)$;
Equation~\eqref{eq:merkle-opening-payload}.\\
\code{5081}&H&recursive opening&$(\bot,k,\bot)$, $1\le k<J_m-1$;
Equation~\eqref{eq:merkle-opening-payload}.\\
\code{5082}&H&final opening&$(\bot,J_m-1,\bot)$;
Equation~\eqref{eq:merkle-opening-payload}.\\
\code{5101}&S&fold challenge&$(\bot,k,i)$; one $\K$ challenge.\\
\code{5102}&S&OOD point&$(a,k,c)$; one $\K$ coordinate, with $c$ increasing.\\
\code{5103}&S&glue challenge&$(\bot,k,i)$ as fixed below; one $\K$ challenge.\\
\code{5104}&S&PoW seed&$(\bot,k,i)$; 16 bytes; request auxiliary payload $C$.\\
\code{5105}&S&query draw&$(\bot,k,a)$; one $\K$ challenge for attempt $a$.\\
\code{5106}&S&query alpha&$(\bot,k,c)$; one $\K$ challenge.\\
\bottomrule
\end{longtable}
\endgroup

The GKR ranges are black-box namespaces, not permission to use
implementation-defined bytes.  The fixed GKR protocol must publish a ledger
with the same columns as this table.  The fixed
\code{3001} event resolves the PCS-level phase-$c$ placement ambiguity; the
mandatory subordinate ledger fixes all remaining GKR messages and squeezes,
and no protocol may reassign the domain code.  Domain code \code{3081} is
reserved and \mustnot{} occur.
The ledger first consumes $\pi_{\rm GKR}$ to derive
$(\bm\omega_\xi,\sigma)$ and then, with no intervening absorb or squeeze,
consumes $\pi_{\rm MLE}$ to derive $(\bm r_h,\theta,\eta)$.  It must require
$\theta=\widetilde{\bm\omega_\xi}(\bm r_h)\eta$ before returning the successor
claim $\widetilde H_{\K}(\bm r_h)=\eta$.  The semantic P5/V5--P6/V6 split does
not create a transcript event.  An absent or incomplete fixed ledger is a hard
reject.

\subsection{Transcript transitions and replay}
\label{sec:wire-transcript}

Model the transcript state as the byte string hashed by an unkeyed BLAKE3
hasher.  It starts empty.  The domain/version value \code{0001} and statement
value \code{0002} are absorbed, in that order, before any proof-dependent
squeeze.
The PIOP-supplied linear-claim statement digest is
\begin{equation}
 d_{\rm claim}:=\operatorname{BLAKE3}_{256}
 \bigl(\operatorname{hex}^{-1}
   (\code{46325a2d5043532f6170702d73746174656d656e742f763100})
 \parallel\mathsf{linear\_claim\_statement\_bytes}\bigr).
 \label{eq:claim-statement-digest}
\end{equation}
For v1 byte compatibility, Equation~\eqref{eq:claim-statement-digest} retains
the legacy ASCII domain literal \code{F2Z-PCS/app-statement/v1} followed by one
NUL byte.  Changing that literal requires a new wire version.  The two initial
payloads are exactly
\begin{equation}
 \begin{aligned}
  \mathsf{DomainV1}:={}&
    \operatorname{hex}^{-1}(\code{46325a5043530000})
    \parallel\LE{16}{1},\\
  \mathsf{StatementV1}:={}&
    d_{\rm claim}[32]
    \parallel\LE{32}{d_r}\parallel\LE{32}{d_c}\parallel\LE{32}{m_f}\\
    &\parallel\operatorname{Enc}_{\K}(g)\parallel\rho[32].
 \end{aligned}
 \label{eq:domain-statement-payloads}
\end{equation}
Thus $\mathsf{DomainV1}$ and $\mathsf{StatementV1}$ have 10 and 92 bytes,
respectively.  The host must make the linear-claim statement byte
string in Equation~\eqref{eq:claim-statement-digest} canonical.  That string
must bind $\operatorname{Vec}_{\Fq}(\bm v)$, $\operatorname{Enc}_{\Fq}(\mu)$,
the configured shapes and coordinate orders of $F$ and $H$, and a
virtualization-mode byte.  The vector encoding is
$\LE{32}{n_h}\parallel\prod_{z=0}^{n_h-1}
\operatorname{Enc}_{\Fq}(v[z])$.  The direct mode binds no matrix and requires
the effective map to be the identity; the virtualized mode also binds the
canonical dimensions and entries of $\bm M$.  Consequently an omitted map
and an explicitly configured nonidentity map cannot share a statement digest.

After the final input-Sumcheck challenge, the verifier derives and absorbs
\begin{equation}
 \mathsf{LinCheckSeamV1}:=
  \LE{32}{d_r}\parallel\LE{32}{d_c}\parallel
  \operatorname{Vec}_{\Fq}(\bm s_h)\parallel
  \operatorname{Enc}_{\Fq}(\kappa)\parallel
  \operatorname{Enc}_{\Fq}(\tau).
 \label{eq:lincheck-seam-payload}
\end{equation}
Here $\operatorname{Vec}_{\Fq}(\bm s_h)$ has its canonical $u_{32}$ count
$m_h$ followed by the $m_h$ canonical coordinates in transcript order.
Equation~\eqref{eq:lincheck-seam-payload} has exactly $44+16m_h$ bytes.  It is
absorbed under derived domain code \code{1003} after the last accepted domain code \code{1101}
challenge and before the domain code \code{2001} record.  The verifier derives
$\kappa=\widetilde{\bm v}(\bm s_h)$,
$a_r=\eq_{\Fq}(r,\bm s_r)$, and
$a'_c=\kappa\eq_{\Fq}(c,\bm s_c)$; none of these dense factor vectors is
serialized in the seam.

Absorption appends exactly Equation~\eqref{eq:event-absorb-v1}.  A squeeze
under base domain code $\tau<\mathtt{0x8000}$, verifier-derived scope
$(j,k,i)$, output length $n$, and auxiliary request payload $a$ first sets
$q:=\LE{32}{n}\parallel a$ and then performs the transition
\begin{equation}
 \begin{aligned}
  q&:=\LE{32}{n}\parallel a,\\
  \tr'&:=\tr\parallel D_{\tau,j,k,i}
    \parallel\LE{64}{\ByteLen(q)}\parallel q,\\
  z&:=\operatorname{BLAKE3\text{-}XOF}(\tr')[0\mathbin{..}n],\\
  \tr''&:=\tr'\parallel
    D_{\tau\mathbin{\mathrm{OR}}\mathtt{0x8000},j,k,i}
    \parallel\LE{64}{n}\parallel z.
 \end{aligned}
 \label{eq:transcript-v1}
\end{equation}
Equation~\eqref{eq:transcript-v1} uses a clone/finalize-XOF of the incremental
hasher after the request bytes; it does not reset or key the hasher.  A $\K$ challenge has
$n=16$ and decodes by Section~\ref{sec:wire-primitives}.  A vector challenge
is a sequence of scalar squeezes, not one long squeeze.  Unless its registry
row says otherwise, the scalars use consecutive $i$ values and the auxiliary
request payload $a$ is empty.

Domain code \code{1101} samples one unbiased $\Fq$ challenge.  For attempts
$a=0,1,\ldots$, apply Equation~\eqref{eq:transcript-v1} with $n=16$ and
auxiliary payload $\LE{32}{a}$.  Interpret the raw result as a little-endian
$u_{128}$, clear its high $28$ bits, and accept the resulting candidate iff it
is less than $q=2^{100}-15$; otherwise continue with the next attempt.  Every
request and raw domain code \code{9101} result remains in the transcript, including
rejected attempts.  Reject before attempt $2^{32}$; the counter must not wrap.

Domain code \code{5103} uses its one-byte glue kind, and domain code \code{5104} uses the PoW
context $C$; domain codes \code{1101}, \code{5103}, and \code{5104} are the only core
squeezes with nonempty auxiliary payloads.

Immediately after domain codes \code{0001} and \code{0002}, the verifier runs exactly
$m_h$ input-Sumcheck rounds.  Starting with
$\sigma_{-1}^{\rm shape}=\mu$, round $i$ reads one payload, requires
$g_i^{\rm shape}(0)+g_i^{\rm shape}(1)=\sigma_{i-1}^{\rm shape}$, absorbs the
payload under domain code \code{1002}, and then samples
$s_{h,i}$ under domain code \code{1101} at the same scope, and sets
$\sigma_i^{\rm shape}=g_i^{\rm shape}(s_{h,i})$.  No event may intervene
between that read and its accepted challenge.  After the final round it sets
$\tau=\sigma_{m_h-1}^{\rm shape}$ and
$\kappa=\widetilde{\bm v}(\bm s_h)$, then immediately absorbs the derived
domain code \code{1003} seam.  The remaining core must discharge
$\kappa\widetilde H_q(\bm s_h)=\tau$; the seam does not assert that the
original $\bm v$ equals the derived equality weights.

On a normal NARG read, the verifier derives the next expected domain and
payload decoder from the ledger, validates the canonical payload, and
immediately applies Equation~\eqref{eq:event-absorb-v1}.  A PoW nonce is the
sole deferred NARG read: the verifier parses its expected canonical payload,
permits no intervening event, checks Equation~\eqref{eq:pow-predicate-v1}, and
absorbs it only on success.  On a Hint read, it derives the next expected
decoder and validates the canonical payload but never absorbs it.  It must
complete the hint predicate before the next squeeze.  In particular, it reads
and absorbs \code{2001}, then reads \code{2081} and requires
$|\bm\mu|=2^{d_c}$, $0\le\mu_c<2^{127}$ and $g^{\mu_c}=\nu_c$ for every $c$,
and $\sum_ca'_c\pi_q(\mu_c)=\tau$ before the first \code{2101} event.  NARG
replay follows verifier-event order, never Rust
structure-field order.  Acceptance requires no deferred record or pending
hint and exact logical and byte exhaustion of both streams and every nested
slice.

After the adjoint, both parties derive the same canonical
$\mathsf{RingDecomp}(\bm a_f)$ and require
$A_f[v,j]=\sum_{\theta\in\Theta}A_\theta[v]B_\theta[j]$.  For each
$\theta\in\Theta$ in order, the verifier reads and absorbs one domain-code
\code{4001} vector and requires $|\bm s_\theta|=128$.  After all such records,
and before domain code \code{4101}, it checks the single anchor equation
\[
 \sum_{\theta\in\Theta}\sum_{v=0}^{127}
   A_\theta[v]s_{\theta,v}=\eta.
\]
It then squeezes the seven coordinates of $r''$ under domain code \code{4101},
derives $(\bm\omega_\theta,\beta_\theta)$ for every term, and squeezes one
domain-code \code{4102} coefficient $\lambda_\theta$ per term in the same
order.  Finally it derives
$B=\sum_\theta\lambda_\theta\bm\omega_\theta$ and
$\beta=\sum_\theta\lambda_\theta\beta_\theta$.  No post-adjoint LinCheck or
coefficient-Sumcheck event exists; domain codes \code{4002} and \code{4103}
are reserved.

\subsection{Grinding and query derivation}
\label{sec:wire-pow}

Let kind byte zero denote a fold grind and kind byte one a query grind.  For
kind $K_{\rm pow}$, level $k$, round $i$, and derived difficulty $d$, define
the exact PoW context
\begin{equation}
 C:=\operatorname{ASCII}(\code{F2Z-PCS/PoW/v1})\parallel
 \LE{8}{K_{\rm pow}}\parallel\LE{32}{k}\parallel
 \LE{32}{i}\parallel\LE{16}{d}.
 \label{eq:pow-context-v1}
\end{equation}
In Equation~\eqref{eq:pow-context-v1}, $i=\bot_{32}$ for a query grind.  The seed is the 16-byte output of
Equation~\eqref{eq:transcript-v1} under domain code \code{5104}, the same scope, and
auxiliary request payload $C$.  After obtaining the seed, the verifier reads
the expected \code{5007} or \code{5008} NARG nonce payload without absorbing it and
requires, for the scheduled positive difficulty $d$,
\begin{equation}
 \LZ\!\left(\operatorname{BLAKE3}_{256}
   (C\parallel\mathsf{seed}_{16}\parallel\LE{64}{n})\right)\ge d,
 \qquad d>0.
 \label{eq:pow-predicate-v1}
\end{equation}
$\LZ$ in Equation~\eqref{eq:pow-predicate-v1} counts zero bits from the
most-significant bit of digest byte zero, then proceeds toward the
least-significant bit and subsequent bytes.  Any satisfying nonce is accepted;
minimality is not checked.  Only after the predicate succeeds does
the verifier absorb the nonce under its verifier-derived domain, before any
dependent squeeze.
Difficulty zero has no predicate invocation, seed squeeze, or nonce record in
this wire format.

The exact scheduled difficulties are
\begin{equation}
 d^{\rm query}_k:=g_k=16,
 \qquad
 d^{\rm fold}_{k,i}:=\max(f_k-i,0).
 \label{eq:pow-difficulties-v1}
\end{equation}
Equation~\eqref{eq:pow-difficulties-v1} is evaluated with ordinary
nonnegative integers.  After the preceding sumcheck message and immediately
before fold challenge $(k,i)$, perform both the domain code \code{5104} seed squeeze
and domain code \code{5007} nonce record if $d^{\rm fold}_{k,i}>0$; if the difficulty
is zero, perform neither event.  For $0\le k<J_m-1$, query event $k$ occurs after the
domain code \code{5003} commitment root for level $k+1$, all $o_{k+1}$ OOD bindings
for that root, and their glue challenges, and immediately before the domain code \code{5105}
draws that open level $k$.  Query event $J_m-1$ instead occurs after the
domain code \code{5006} final $y_r$, with no new root or OOD binding, and immediately
before the final-level draws.  There is exactly one query nonce at every level.
Consequently the exact number of nonce records is
\begin{equation}
 N_{\rm pow}=J_m+\sum_{k=0}^{J_m-1}\min(k_k,f_k).
 \label{eq:pow-count-v1}
\end{equation}
Missing, extra, reordered, or trailing nonces violate
Equation~\eqref{eq:pow-count-v1} and reject.

For level $k$, each \code{5105} event yields $r\in\K$ and proposes
\begin{equation}
 q:=r_{\lo}\bmod 2^{h_k+r_k}.
 \label{eq:query-position-v1}
\end{equation}
The verifier repeats Equation~\eqref{eq:query-position-v1} until $Q_k$
distinct positions exist, discarding duplicates, and then sorts the positions
in increasing order.  It then derives the domain code \code{5106} alphas before
reading and checking the opening.

\subsection{Recursive-opening replay ledger}
\label{sec:wire-recursive-ledger}

This subsection fixes the complete replay order for the recursive opening.
The four canonical message kinds and scopes for domain code \code{5004} are
\begin{equation}
 \begin{array}{ccll}
  \mathtt{0x00}&\mathsf{Start}&(\bot,0,\bot)&\text{exactly once},\\
  \mathtt{0x01}&\mathsf{FoldNext}&(\bot,k,j)&0\le j<k_k,\\
  \mathtt{0x02}&\mathsf{OodIntro}&(a,k,\bot)&1\le k<J_m, 0\le a<o_k,\\
  \mathtt{0x03}&\mathsf{QueryIntro}&(\bot,k,\bot)&0\le k<J_m-1.
 \end{array}
 \label{eq:sumcheck-message-kinds}
\end{equation}
Equation~\eqref{eq:sumcheck-message-kinds} is the complete kind registry for
domain code \code{5004}.  Each payload is the one-byte kind followed by the
two canonical $\K$
elements in the registry.  Any other kind/scope pair rejects.

The outer $\mathsf{VerifyF2Z}$ caller uses the deterministic coefficient
function $B$ derived locally in P8/V8, with domain size
$2^{m_p}$; $B$ is never densely serialized.  The caller alone absorbs domain code
\code{5001} with payload
$\LE{32}{m_p}\parallel\operatorname{Enc}_{\K}(\beta)$ and then invokes
$\mathsf{RecVerify}$.  That helper starts after domain code \code{5001}: its
first event reads domain code \code{5002}, requires the root to equal $\rho$, absorbs
it, and then reads and absorbs the $\mathsf{Start}$ domain code \code{5004} message.

For each $k=0,\ldots,J_m-1$ in increasing order, it then performs the
following ledger without interleaving another event.
\begin{enumerate}
 \item For $j=0,\ldots,k_k-1$, perform the optional fold grind exactly as in
 Section~\ref{sec:wire-pow}, squeeze domain code \code{5101} at
 $(\bot,k,j)$, then read, validate, and absorb the $\mathsf{FoldNext}$
 domain code \code{5004} message at the same scope.

 \item If $k<J_m-1$, read and absorb the domain code \code{5003} root at
 $(\bot,k+1,\bot)$.  For each $a=0,\ldots,o_{k+1}-1$, squeeze the
 $n_{k+1}$ domain code \code{5102} coordinates at scopes $(a,k+1,c)$ for
 $c=0,\ldots,n_{k+1}-1$; read and absorb domain code \code{5005} at
 $(a,k+1,\bot)$; read and absorb the $\mathsf{OodIntro}$ domain code \code{5004}
 message at that scope; and squeeze domain code \code{5103} at
 $(\bot,k+1,a)$ with the one-byte auxiliary kind \code{00}.

 If instead $k=J_m-1$, read and absorb domain code \code{5006} at
 $(\bot,k,\bot)$, requiring exactly $2^{h_k}$ canonical $\K$ elements.

 \item Perform the query grind at $(\bot,k,\bot)$ with difficulty 16 and
 absorb its validated domain code \code{5008} nonce.  Starting with attempt $a=0$,
 squeeze domain code \code{5105} at $(\bot,k,a)$ and increment $a$ after every
 squeeze, including a duplicate, until $Q_k$ distinct positions have been
 obtained.  Reject rather than wrap if the next attempt would exceed
 $2^{32}-1$.  Sort the positions, then squeeze exactly
 $\lceil\log_2Q_k\rceil$ domain code \code{5106} alphas at $(\bot,k,c)$ for
 consecutive $c$.

 Read exactly one opening Hint using domain code \code{5080} for $k=0$,
 \code{5081} for $0<k<J_m-1$, or \code{5082} for $k=J_m-1$.  Check
 Equations~\eqref{eq:merkle-opening-payload}--\eqref{eq:merkle-opening-length}
 and the root before any later event.  For $k<J_m-1$, next read and absorb
 the $\mathsf{QueryIntro}$ domain code \code{5004} message.  At every level,
 including the final level where no such message exists, squeeze the query
 glue domain code \code{5103} at $(\bot,k,o_k)$ with one-byte auxiliary kind
 \code{01}.  At the final level this challenge is used directly in the
 residual check.
\end{enumerate}

The fixed record/event cardinalities are therefore
\begin{equation}
 \begin{aligned}
  \#5002&=1,& \#5003&=J_m-1,& \#5005&=\sum_{k=1}^{J_m-1}o_k,\\
  \#5004&=1+\sum_{k=0}^{J_m-1}k_k
     +\sum_{k=1}^{J_m-1}o_k+(J_m-1),&
  \#5006&=1,& \#508{*}&=J_m,\\
  \#5101&=\sum_{k=0}^{J_m-1}k_k,&
  \#5102&=\sum_{k=1}^{J_m-1}o_kn_k,&
  \#5103&=J_m+\sum_{k=1}^{J_m-1}o_k,\\
  \#5106&=\sum_{k=0}^{J_m-1}\lceil\log_2Q_k\rceil.
 \end{aligned}
 \label{eq:recursive-event-counts}
\end{equation}
The cardinalities in Equation~\eqref{eq:recursive-event-counts} are exact.
Here $\#508{*}$ is one opening Hint per level.  The number of domain-
\code{5105} events is the
transcript-derived number of attempts, all of which appear in the audit log.
Nonce counts remain Equation~\eqref{eq:pow-count-v1}.  A different order,
scope, kind, cardinality, or unread field rejects.

\subsection{Parser, implementation, and conformance requirements}
\label{sec:wire-conformance}

A v1 implementation \must{}:
\begin{enumerate}
 \item validate magic, version, flags, all public-parameter-, statement-, and
 transcript-derived bounds, and all
reserved fields before allocating or doing proof work;
 \item use checked arithmetic and checked conversion for every length and
 index, reject noncanonical field values, and return a decoding error rather
 than panic on adversarial bytes;
 \item encode the NARG stream in the exact absorb order and the Hint stream in
 its exact check order, with independent cursors;
 \item replace opaque recursive-proof serialization with the field records and
 exact order in Sections~\ref{sec:wire-events}
 and~\ref{sec:wire-recursive-ledger}; and
 \item reject unless every outer, subordinate, recursive, row, and Merkle-path
 slice is exactly exhausted.
\end{enumerate}

Conformance fixtures \must{} include canonical primitive, domain, and payload
hex, one
complete $N/H/\HostEncode$ proof, the ordered transcript audit log, every
squeezed challenge, PoW seeds/nonces and query positions, and all four lengths
from Equation~\eqref{eq:wire-lengths-v1}.  Negative fixtures \must{} cover an
unknown version, wrong endian or channel, record
reordering, truncation and suffix bytes, noncanonical scalars, count overflow,
a wrong input-Sumcheck round polynomial, wrong initial target, reordered
\code{1002/1101} events, a changed domain code \code{1003} seam, a failed or late
$\mu/\nu$ hint, altered PoW scope/difficulty/nonce, and missing or extra
recursive fields.  Fixtures must also cover at least one rejected
domain code \code{1101} sampling attempt and its transcript continuation.
Ring-switch fixtures must additionally cover a malformed or reordered
\code{4001} term, a failed anchor equation, and reordered or missing
\code{4101/4102} challenges.  Native, recursive, and independent reference
verifiers must produce byte-identical transcript audit logs.

The repository file \code{wire-v1-vectors.txt} is the normative initial
known-answer set for the primitive container/domain/payload codec,
Merkle row/tree/multiproof codec, and PoW predicate at difficulties 1, 7, 8,
and 16.  A conforming implementation must reproduce it exactly and add a
complete proof/transcript
fixture before claiming end-to-end v1 interoperability.

\endgroup


\end{document}
